\documentclass[11pt]{amsart}

\usepackage[T1]{fontenc}
\usepackage{amsmath}
\usepackage{amssymb}

\DeclareMathOperator{\lcm}{lcm}
\DeclareMathOperator{\num}{num}
\DeclareMathOperator{\rad}{rad}
\newcommand{\Z}{\mathbb{Z}}

\theoremstyle{plain}
\newtheorem{theorem}{Theorem}
\renewcommand{\thetheorem}{\Alph{theorem}}
\newtheorem{lemma}{Lemma}
\newtheorem{proposition}{Proposition}
\newtheorem{corollary}{Corollary}
\newtheorem{conjecture}{Conjecture}
\newtheorem{hypothesis}{Hypothesis}
\theoremstyle{remark}
\newtheorem{remark}{Remark}

\title{The Ap\'ery gcd: a density-one theorem and the pointwise endpoint}
\author{X. Huang}
\date{}

\begin{document}
\raggedbottom

\begin{abstract}
Let $a_n$ and $b_n$ be the two standard solutions of the Ap\'ery
recurrence for $\zeta(3)$.  Put
\[
 d_n=\lcm(1,\ldots,n)^3,
 \qquad G_n=\gcd(d_na_n,d_nb_n).
\]
We prove that for every $\varepsilon>0$ the set of $n\leq X$ for which
$\log G_n>\varepsilon n$ has cardinality $o(X)$.
More precisely, on every dyadic shell $(N,2N]$ the average of
$\log G_n$ is $O(N^{2/3}\log N)$.  For the pointwise problem we
identify the exact radical middle prime channel, prove a reflection
functional equation, and construct fixed gap-eliminant polynomials
$U_g$.  In the only delicate range we lower the pointwise problem to a
depth-$2$ pair gate and then to a single-arc exact-collision conjecture.
We prove complementary cross-arc and structured-resonance results, give
two algebraic descriptions of the zero sets, and record barrier
certificates showing why horizontal equidistribution does not resolve
the remaining scale-zero collision problem.
\end{abstract}

\maketitle

\section{Statement}

Ap\'ery's construction \cite{Apery1979} gives two $\zeta(3)$ solutions
$b_n$ and $a_n$ of
\[
 (n+1)^3u_{n+1}
 = (2n+1)(17n^2+17n+5)u_n-n^3u_{n-1},
\]
with $(b_0,b_1)=(1,5)$, where the $b_n$ are integers, and
$(a_0,a_1)=(0,6)$, where the $a_n$ are rational and
\[
 v_p(a_n)\geq -3\lfloor\log_p n\rfloor.
\]
Set
\[
 d_n=\lcm(1,\ldots,n)^3,
 \qquad
 G_n=\gcd(d_na_n,d_nb_n).
\]

\begin{theorem}\label{thm:density-one}
For every $\varepsilon>0$,
\[
 \#\{n\leq X:\log G_n>\varepsilon n\}=o(X).
\]
Quantitatively, on every dyadic shell,
\[
 \frac1N\sum_{N<n\leq2N}\log G_n=O(N^{2/3}\log N).
\]
\end{theorem}

Write
\[
 P(t)=(2t+1)(17t^2+17t+5)
\]
throughout.  Unless stated otherwise, $p$ always denotes a prime at least
$7$.  The primes $2,3,5$ contribute $O(\log^2 n)$ to $\log G_n$, since
\[
 v_p(G_n)\leq v_p(d_n)+O(\log_p n)=O(\log n),
\]
and they are absorbed into the small-prime estimate.

\begin{remark}[Provenance and status]
This proof chain was assembled on July 19, 2026, by extracting and
streamlining \texttt{Q3.2\_paper\_draft.md}.  Every lemma used below is
included, and a summary of the accompanying exact verification appears in
the appendix.
\end{remark}

\section{Preliminaries}

\subsection{Wronskian}

Define
\[
 W_n:=a_{n+1}b_n-a_nb_{n+1}.
\]
Then
\[
 W_n=\frac{6}{(n+1)^3}.
\]
Indeed, $W_0=6$, and the recurrence gives
$W_n=\frac{n^3}{(n+1)^3}W_{n-1}$.  In particular, $W_n\neq0$, and for
$\sqrt n<p\leq n$, with $p\geq7$,
\[
 v_p(G_n)\leq6.
\]

\begin{proof}
The integer $G_n$ divides both $d_na_n$ and $d_nb_n$, hence it divides the
integer combination
\begin{align*}
 &d_{n+1}b_{n+1}(d_na_n)
 -d_{n+1}a_{n+1}(d_nb_n)\\
 &\qquad=-d_nd_{n+1}W_n\\
 &\qquad=-\frac{6d_nd_{n+1}}{(n+1)^3}.
\end{align*}
For $\sqrt n<p\leq n$, one has
$v_p(d_n),v_p(d_{n+1})\leq3$ and $p\nmid6$.  Therefore
\[
 v_p(G_n)\leq6+v_p(6)-3v_p(n+1)\leq6.
\]
\end{proof}

\begin{remark}[Audit correction]
Any absolute constant depth suffices for the averaging argument.  The
bound $6$ above replaces an earlier unproved bound $4$ identified during
the hostile audit.
\end{remark}

\subsection{Cleared continuants}

Let
\[
 N_0=0,
 \qquad
 N_1=1,
 \qquad
 N_{r+1}(x)=P(x+r)N_r(x)-(x+r)^6N_{r-1}(x)\in\Z[x].
\]
Then $\deg N_r=3(r-1)$ for $r\geq1$.  Every solution of the local cleared recurrence
at base $x$ satisfies
\begin{equation}
 \begin{aligned}
 Q_h(x)y_h
 &=N_h(x)y_1-(x+1)^3N_{h-1}(x+1)y_0,\\
 Q_h(x)&=\prod_{i=2}^h(x+i)^3.
 \end{aligned}
 \tag{1.2}\label{eq:solution-formula}
\end{equation}

\subsection{Lucas congruence}

If $n=\sum_i d_ip^i$ is the base-$p$ expansion of $n$, then
\[
 b_n\equiv\prod_i b_{d_i}\pmod p.
\]
This is Gessel's Lucas congruence \cite{Gessel1982}; related divisibility
and Lucas phenomena for Ap\'ery-like sequences are developed in
\cite{MalikStraub2016,LucaShparlinski2008}.

\section{Five continuant lemmas}

\begin{lemma}[No consecutive zeros; both solutions]\label{lem:no-consecutive}
For $p\geq7$ and $0\leq r\leq p-2$, it is not the case that both
$b_r\equiv0\pmod p$ and $b_{r+1}\equiv0\pmod p$.  For
$1\leq q\leq p-2$, it is not the case that both
$a_q\equiv0\pmod p$ and $a_{q+1}\equiv0\pmod p$; here $a_q$ is
$p$-integral because $q<p$.
\end{lemma}

\begin{proof}
Two consecutive zeros back-propagate through the recurrence, because all
coefficients $m^3$ with $m<p$ are invertible.  For $b$ this reaches
$b_0=1$, a contradiction.  For $a$ it reaches $a_1=6$, and $p\nmid6$.
\end{proof}

\begin{lemma}[Evaluation identities]\label{lem:evaluation}
For every $r\geq1$,
\begin{enumerate}
\item[(i)] $N_r(-1)=((r-1)!)^3b_{r-1}$;
\item[(ii)] $N_r(-x-r)=(-1)^{r-1}N_r(x-1)$, and consequently
\[
 N_r(-r)=(-1)^{r-1}((r-1)!)^3b_{r-1}.
\]
\end{enumerate}
\end{lemma}

\begin{proof}
For the first identity, $v_r:=N_r(-1)$ satisfies $v_1=1$,
$v_2=P(0)=5$, and
\[
 v_{r+1}=P(r-1)v_r-(r-1)^6v_{r-1}.
\]
Substituting $v_r=((r-1)!)^3b_{r-1}$ reduces this to the Ap\'ery
recurrence at $n=r-1$.

For the second identity, under $x\mapsto-x-r$ each diagonal entry of the
defining tridiagonal matrix transforms by $P(-1-t)=-P(t)$, and each
off-diagonal cube is negated, in reversed order.  Simultaneous
row-and-column reversal preserves the determinant, while the negation
contributes $(-1)^{r-1}$ because the matrix has size $r-1$.
\end{proof}

\begin{lemma}[Continuant lemma]\label{lem:continuant}
For $p\geq7$ and $1\leq h<p$, one has
$N_h\not\equiv0\pmod p$.  For $p=5$ the same assertion follows by a
direct check for $h<5$.
\end{lemma}

\begin{proof}
Suppose that $N_h\equiv0\pmod p$.  Lemma~\ref{lem:evaluation}(i) gives
$b_{h-1}\equiv0\pmod p$.  The defining recurrence, evaluated at
$x=1-h$, gives
\[
 N_h(1-h)=P(0)N_{h-1}(1-h)=5N_{h-1}(1-h).
\]
Since $p\neq5$, it follows that
$N_{h-1}(1-h)=N_{h-1}(-(h-1))\equiv0\pmod p$.  By
Lemma~\ref{lem:evaluation}(ii), $b_{h-2}\equiv0\pmod p$, contrary to
Lemma~\ref{lem:no-consecutive}.
\end{proof}

\begin{lemma}[Pair certificate]\label{lem:pair-certificate}
Let $y$ be any solution of the Ap\'ery recurrence modulo $p$ on a window
$[m,m+h]$, where $m+h<p$.  If
\[
 y_m\equiv y_{m+h}\equiv0\pmod p
 \qquad\text{and}\qquad
 y_{m+1}\not\equiv0\pmod p,
\]
then $N_h(m)\equiv0\pmod p$.
\end{lemma}

\begin{proof}
Formula~\eqref{eq:solution-formula} at base $x=m$, with $y_0=0$, gives
\[
 Q_h(m)y_{m+h}=N_h(m)y_{m+1}.
\]
All factors $m+i$ for $2\leq i\leq h$ lie in
$(m,m+h]\subseteq(0,p)$, so $Q_h(m)$ is a unit.  The two hypotheses on
$y_{m+h}$ and $y_{m+1}$ force $N_h(m)\equiv0\pmod p$.
\end{proof}

\begin{lemma}[Zero-fiber bounds]\label{lem:zero-fiber}
For $p\geq7$,
\[
 \begin{aligned}
 \mathcal Z_b(p)&:=\{0\leq r<p:b_r\equiv0\pmod p\},\\
 Z_b(p)&:=|\mathcal Z_b(p)|\leq3p^{2/3}+O(1),
 \end{aligned}
\]
and
\[
 \begin{aligned}
 \mathcal Z_a(p)&:=\{1\leq q<p:a_q\equiv0\pmod p\},\\
 Z_a(p)&:=|\mathcal Z_a(p)|\leq3p^{2/3}+O(1).
 \end{aligned}
\]
\end{lemma}

\begin{proof}
Let $r_1<\cdots<r_z$ be the zeros and set
$h_i=r_{i+1}-r_i$, so that $\sum_i h_i<p$.  By
Lemmas~\ref{lem:no-consecutive} and~\ref{lem:pair-certificate}, every
consecutive pair with $h_i\leq K$ makes $r_i$ a root of
$\overline N_{h_i}\neq0$; the nonvanishing follows from
Lemma~\ref{lem:continuant}.  For a fixed $h$ the bases are distinct.
Since $\deg N_h=3(h-1)$, at most
\[
 \sum_{h\leq K}3(h-1)<\frac32K^2
\]
pairs have gap at most $K$, while fewer than $p/K$ have gap greater than
$K$.  Taking $K=p^{1/3}$ proves the bound for $Z_b(p)$.  The $a$-case is
identical; the zero $a_0=0$ is excluded from the count, and its windows
start at $q\geq1$.
\end{proof}

\section{The support law}

Write $\widetilde a_n=p^3a_n$ where relevant; this quantity is
$p$-integral for $n<p^2$.  To make the binomial notation used below
explicit, set
\[
 H_N^{(3)}=\sum_{j=1}^N\frac1{j^3},
 \qquad
 T(N,K)=\binom NK^2\binom{N+K}K^2,
\]
and
\[
 E(N,K)=H_N^{(3)}+
 \sum_{j=1}^K
 \frac{(-1)^{j-1}}
 {2j^3\binom Nj\binom{N+j}j}.
\]
The second Ap\'ery solution has the standard representation
\[
 a_N=\sum_{K=0}^N T(N,K)E(N,K).
\]
For background on congruences for these and closely related Ap\'ery sums,
see \cite{Mortenson2005,Sun2012,WangSun2024}.

\begin{lemma}[The congruence $\alpha_p\equiv6$]\label{lem:alpha-six}
For $p\geq7$,
\[
 p^3a_p\equiv6\pmod p.
\]
\end{lemma}

\begin{proof}
In the binomial sum
$a_p=\sum_{k=0}^pT(p,k)E(p,k)$, the $p^3$-pole survives only at
$k=0$ and $k=p$.  For $1\leq k\leq p-1$, one has
$v_p(\binom pk^2)\geq2$.  After multiplication by $p^3$, the
$H_p^{(3)}$ part therefore has valuation at least $2$, while the inner
sum, whose $p$-adic valuation is at least $-1$, has valuation at least
$4$.  Thus all these terms vanish modulo $p$.  The $k=0$ term contributes
\[
 p^3H_p^{(3)}=p^3H_{p-1}^{(3)}+1\equiv1\pmod p.
\]
The $k=p$ term contributes
\[
 \binom{2p}p^2\left(1+\frac14\right)
 \equiv4\cdot\frac54=5\pmod p.
\]
Here Lucas gives $\binom{2p}p\equiv2\pmod p$; the $j=p$ inner term
gives
\[
 \frac{p^3}{2p^3\binom{2p}p}\equiv\frac14\pmod p,
\]
and the inner terms with $j<p$ have pole order at most $1$.  The total is
$1+5=6$.
\end{proof}

\begin{lemma}[Boundary lemma]\label{lem:boundary}
For $1\leq q<p$ and $p\geq7$,
\[
 p^3a_{qp}\equiv a_q\pmod p.
\]
\end{lemma}

\begin{proof}[Kummer bookkeeping]
In the sum for $a_{qp}$, $T(qp,K)$ is a $p$-unit if and only if
$K=\ell p$ with $q+\ell<p$.  Indeed, by Kummer, any $s>0$ in
$K=\ell p+s$ forces borrows.  Together with the carry indicator of
$\binom{qp+K}K$, this gives total valuation at least $2$.  The bound
$v_p(E(qp,K))\geq-3$ shows that non-unit terms die after multiplication
by $p^3$; termwise, the inner terms with $p\nmid j$ have total valuation
at least $3$, and the terms $j=tp$ have total valuation at least $1$.

On the surviving terms, Lucas gives
$T(qp,\ell p)\equiv T(q,\ell)\pmod p$; moreover,
\[
 p^3H_{qp}^{(3)}\equiv H_q^{(3)}\pmod p,
\]
because only the multiples of $p$ survive, and the inner sum reduces
index by index with $(-1)^{tp-1}=(-1)^{t-1}$.  This yields
\[
 \sum_{\ell:\,q+\ell<p}T(q,\ell)E(q,\ell).
\]
Restoring the terms with $q+\ell\geq p$ costs nothing: there
$v_p(T(q,\ell))=2$ from one carry, squared, while
$v_p(E(q,\ell))\geq-1$.
\end{proof}

\begin{lemma}[Block law in the $r$-direction]\label{lem:block-law}
For $0\leq r<p$ and $1\leq q<p$,
\[
 \widetilde a_{qp+r}\equiv\widetilde a_{qp}b_r\pmod p.
\]
\end{lemma}

\begin{proof}
Multiply the recurrence at $n=qp+r$ by $p^3$.  For
$1\leq r\leq p-2$, the coefficients reduce to those in the
$b$-recurrence.  At the entry $r=0$, the backward term carries
$(qp)^3=q^3p^3$ against the $p$-integral number
$\widetilde a_{qp-1}$, and hence vanishes modulo $p$; also
$(qp+1)^3\equiv1\pmod p$, so
\[
 \widetilde a_{qp+1}\equiv P(0)\widetilde a_{qp}
 =5\widetilde a_{qp}\pmod p.
\]
Thus both sides satisfy the same nondegenerate recurrence with initial
data
\[
 \widetilde a_{qp}(1,5)=\widetilde a_{qp}(b_0,b_1).
\]
\end{proof}

\begin{theorem}[Support law]\label{thm:support-law}
Let $p\geq7$ and $\sqrt n<p\leq n$.  Write $n=qp+r$, where
$1\leq q<p$ and $0\leq r<p$.  Then
\[
 p\mid G_n
 \quad\Longleftrightarrow\quad
 a_qb_r\equiv0\pmod p.
\]
\end{theorem}

\begin{proof}
In this range $v_p(d_n)=3$, and $v_p(b_n)\geq0$.  Hence
\[
 p\mid G_n
 \quad\Longleftrightarrow\quad
 v_p(a_n)\geq-2
 \quad\Longleftrightarrow\quad
 \widetilde a_n=p^3a_n\equiv0\pmod p.
\]
By Lemmas~\ref{lem:boundary} and~\ref{lem:block-law},
\[
 \widetilde a_n
 \equiv\widetilde a_{qp}b_r
 \equiv a_qb_r\pmod p.
\]
\end{proof}

\begin{remark}[Historical correction]
The earlier one-channel claim
``$p\mid G_n\Rightarrow b_r\equiv0\pmod p$'' is false outside
$n/2<p<n$: the hostile audit found the counterexamples
$13\mid G_{26}$ and $13\mid G_{27}$.  In the top range $q=1$, however,
$a_1=6$ is a unit, and the old claim is recovered.
\end{remark}

\section{Proof of Theorem A}\label{sec:proof-density}

\begin{proof}[Proof of Theorem A]
Fix the shell $(N,2N]$.  For $n$ in this shell, abbreviate
$v_p=v_p(G_n)$ and write
\[
 \log G_n
 =\underbrace{\sum_{p\leq\sqrt n}v_p\log p}_{S_1(n)}
 +\underbrace{\sum_{\sqrt n<p\leq n}v_p\log p}_{S_2(n)}.
\]

\paragraph{Small primes.}
For every $p\leq\sqrt n$,
\[
 v_p(G_n)\leq v_p(d_n)+\max(0,v_p(a_n))=O(\log_p n).
\]
Consequently,
\[
 S_1(n)=O(\pi(\sqrt n)\log n)=O(\sqrt n).
\]

\paragraph{Large primes.}
By Theorem~\ref{thm:support-law}, a prime $p$ is bad if and only if
$b_r\equiv0\pmod p$ in the lower channel or
$a_q\equiv0\pmod p$ in the companion channel.  By the Wronskian depth
estimate, each bad prime costs at most $6\log(2N)$.

For the lower channel, fix $p$.  The shell integers with
$n\bmod p$ in the zero fiber counted by $Z_b(p)$ number at most
\[
 Z_b(p)\left(\frac Np+O(1)\right).
\]
The endpoint term is $O(Z_b(p))$ for each prime.  Lemma~\ref{lem:zero-fiber}
therefore gives
\begin{align*}
 \sum_{p\leq2N}Z_b(p)\left(\frac Np+O(1)\right)
 &\ll N\sum_{p\leq2N}p^{-1/3}
     +\sum_{p\leq2N}p^{2/3}\\
 &\ll N^{5/3}.
\end{align*}

For the companion channel, fix $q\geq1$.  The congruence
$a_q\equiv0\pmod p$ means that $p$ divides the fixed nonzero integer
$\num(a_q)$.  This integer is nonzero because $a_q>0$ for $q\geq1$, and
\[
 \log|\num(a_q)|=O(q).
\]
It therefore has at most $O(q/\log N)$ prime divisors exceeding
$\sqrt N$.  Each such pair $(q,p)$ poisons at most
$p\leq2N/q$ shell integers, namely those for which
$\lfloor n/p\rfloor=q$.  The total number of incidences is
\[
 \sum_{q\leq2\sqrt{2N}}
 \frac{2N}{q}\,O\left(\frac q{\log N}\right)
 =O\left(\frac{N^{3/2}}{\log N}\right).
\]

\paragraph{Assembly.}
Summing both channels over the shell and multiplying by the depth cost
gives
\begin{align*}
 \sum_{N<n\leq2N}S_2(n)
 &\ll\left(N^{5/3}+\frac{N^{3/2}}{\log N}\right)\log N\\
 &\ll N^{5/3}\log N.
\end{align*}
Together with the small-prime estimate, this shows that the shell average
of $\log G_n$ is $O(N^{2/3}\log N)$.  Markov's inequality on the shell
now yields
\begin{align*}
 \#\{n\in(N,2N]:\log G_n>\varepsilon n\}
 &\leq\#\{n\in(N,2N]:\log G_n>\varepsilon N\}\\
 &\ll_{\varepsilon}N^{2/3}\log N.
\end{align*}
Summing over dyadic shells up to $X$ gives $o(X)$.
\end{proof}

\section{The pointwise reduction}\label{sec:pointwise}

Put
\[
 b_m=\sum_{k=0}^m\binom mk^2\binom{m+k}k^2,
 \qquad
 Z_p=\{0\leq r<p:b_r\equiv0\pmod p\},
\]
and write $e(t)=\exp(2\pi i t)$ and $\|t\|=\operatorname{dist}(t,\Z)$.

\subsection{Reflection}

\begin{lemma}[Reflection functional equation]\label{lem:reflection}
For every prime $p$ and every $0\leq n\leq p-1$,
\[
 b_{p-1-n}\equiv b_n\pmod p.
\]
\end{lemma}

\begin{proof}
For $0\leq k<p$, the product formula for binomial coefficients gives
\[
 \binom{p-1-m}{k}
 \equiv(-1)^k\binom{m+k}{k}\pmod p.
\]
Apply this first with $m=n$ and then with $m=n-k$.  Thus
\[
 \begin{aligned}
 \binom{p-1-n}{k}&\equiv(-1)^k\binom{n+k}{k},\\
 \binom{p-1-n+k}{k}&\equiv(-1)^k\binom nk\pmod p.
 \end{aligned}
\]
Both Ap\'ery sums may be extended to $0\leq k<p$: outside the original
summation range the corresponding binomial factor is zero modulo $p$.
The two displayed congruences therefore identify their summands term by
term, after squaring, and summation proves the result.
\end{proof}

This direct proof is Lemma~6.2 of Malik--Straub
\cite{MalikStraub2016}; see also the broader congruence work of Beukers
and Coster \cite{Beukers1985,Coster1988}.  The identity was also
checked with no failure for all primes $p\leq12809$; the termwise
congruence was checked for $p\leq809$.

\subsection{The common lower channel}

Let $\sqrt n<p\leq n$ and write $n=qp+r$, with
$q=\lfloor n/p\rfloor$ and $0\leq r<p$.  The support law of
Theorem~\ref{thm:support-law} and the Gessel--Coster Lucas congruence give
two different two-channel decompositions:
\begin{align}
 p\mid G_n
 &\quad\Longleftrightarrow\quad a_qb_r\equiv0\pmod p,
 \label{eq:G-two-channel}\\
 p\mid b_n
 &\quad\Longleftrightarrow\quad b_qb_r\equiv0\pmod p.
 \label{eq:b-two-channel}
\end{align}
The second equivalence follows from
$b_{qp+r}\equiv b_qb_r\pmod p$ \cite{Gessel1982,Coster1988}.
The two decompositions share exactly the condition $r\in Z_p$.

\begin{lemma}[Both digit channels are pointwise-free]
\label{lem:digit-channels}
For every $n$,
\begin{align*}
 \#\{p\in(\sqrt n,n]:a_{\lfloor n/p\rfloor}\equiv0\pmod p\}
 &\ll n^{2/3},\\
 \#\{p\in(\sqrt n,n]:\lfloor n/p\rfloor\in Z_p\}
 &\ll n^{2/3}.
\end{align*}
\end{lemma}

\begin{proof}
For fixed $q$, every prime under consideration lies in
\[
 I_q(n)=\left(\frac n{q+1},\frac nq\right],
 \qquad |I_q(n)|<\frac n{q^2}.
\]
In the first line it divides the fixed nonzero integer $\num(a_q)$; in
the second it divides the fixed positive integer $b_q$.  The standard
growth estimates $\log|\num(a_q)|,\log b_q=O(q)$ show that either integer
has $O(q/\log n)$ prime factors exceeding $\sqrt n$.  Hence each
$q\leq\sqrt n$ contributes at most
\[
 \min\left(\frac n{q^2}+1,\frac{Cq}{\log n}\right).
\]
Splitting the sum at $q=n^{1/3}$ proves both assertions.
\end{proof}

Define the radical lower-channel mass
\[
 R_b(n)=\sum_{\substack{\sqrt n<p\leq n\\n\bmod p\in Z_p}}\log p
\]
and the radical middle part of $b_n$ by
\[
 M_{\mathrm{rad}}(n)
 =\sum_{\substack{\sqrt n<p\leq n\\p\mid b_n}}\log p.
\]

\begin{theorem}[Exact pointwise endpoint]\label{thm:pointwise-reduction}
As $n\to\infty$,
\begin{align}
 R_b(n)
 &\leq\log G_n,
 \label{eq:lower-Rb}\\
 \log G_n
 &\leq6R_b(n)+O(n^{2/3}\log n+\sqrt n),
 \label{eq:upper-Rb}\\
 M_{\mathrm{rad}}(n)
 &=R_b(n)+O(n^{2/3}\log n).
 \label{eq:middle-radical}
\end{align}
Consequently,
\[
 \log G_n=o(n)\quad\text{for every $n$}
 \quad\Longleftrightarrow\quad
 M_{\mathrm{rad}}(n)=o(n)\quad\text{for every $n$}.
\]
\end{theorem}

\begin{proof}
If $n\bmod p\in Z_p$, then \eqref{eq:G-two-channel} gives $p\mid G_n$,
which proves \eqref{eq:lower-Rb}.  Conversely, every prime in the large
range dividing $G_n$ belongs either to the common lower channel or to the
first digit channel in Lemma~\ref{lem:digit-channels}; each has depth at
most $6$.  The small-prime estimate in
Section~\ref{sec:proof-density} is $O(\sqrt n)$, proving
\eqref{eq:upper-Rb}.  Finally, \eqref{eq:b-two-channel} and the second
line of Lemma~\ref{lem:digit-channels} give
\eqref{eq:middle-radical}.
\end{proof}

\begin{conjecture}[Full-multiplicity middle part]
The stronger full-multiplicity statement requested in the original
program is
\begin{equation}
 \log\prod_{\substack{\sqrt n<p\leq n\\p\mid b_n}}
 p^{v_p(b_n)}=o(n).
 \tag{T}\label{eq:T}
\end{equation}
It implies the pointwise theorem, but the converse is not proved: the gcd
sees each middle prime only to depth at most $6$.  Thus
Theorem~\ref{thm:pointwise-reduction}, not \eqref{eq:T} with full
multiplicity, is the unconditional equivalence.
\end{conjecture}

\begin{remark}[Rhin--Viola correction]
The Rhin--Viola irrationality measure \cite{RhinViola2001} gives
\[
 \limsup_{n\to\infty}\frac{\log G_n}{n}
 \leq 3+\log(17+12\sqrt2)
 -\frac{2\log(17+12\sqrt2)}{5.51389\ldots}
 =5.24673\ldots.
\]
The earlier value, $2.24672\ldots$, omitted the contribution
$\log d_n/n\to3$ and is retracted.
\end{remark}

\subsection{The in-range window}

For a dyadic parameter $X$, put
\[
 K_X(y)=\#\{p\in(X,2X]:y\bmod p\in Z_p\}.
\]

\begin{remark}[The required range]\label{rem:required-range}
Fix $A>0$.  A full block $(X,2X]$ is delicate only when
\[
 (X,2X]\subset\left(\frac{n}{(\log n)^A},n\right].
\]
After increasing $A$ harmlessly, all such blocks are covered by the
working window
\[
 \mathcal J_X=\{y\in\Z:2X\leq y\leq X L\},
 \qquad L=(\log X)^A.
 \tag{6.1}\label{eq:in-range-window}
\]
If $X\leq n/(\log n)^A$, the trivial estimate
$K_X(n)\leq\pi(2X)$ closes the block after multiplication by its
logarithmic weight.  Thus no estimate at $n=O(1)$ is used.
\end{remark}

\begin{lemma}[Intercept vacuity]\label{lem:intercept-vacuity}
Suppose $p,p'\in(X,2X]$ are distinct and two incidences lie on the
affine line
\[
 r=up+v,\qquad r'=up'+v.
\]
If their unwrapped phases are $X^{-1}$-close at $n$, then
\[
 |p-p'|\leq\frac{pp'}{X|n-v|}
 \leq\frac{4X}{|n-v|}.
 \tag{6.2}\label{eq:intercept-vacuity}
\]
Consequently, for fixed $B$, if $|v|\ll(\log X)^B$ and $n\geq2X$,
then $|p-p'|\leq2$ for all sufficiently large $X$.  A line carrying
$T$ incidences contributes only $O(T)$ close pairs, independently of
any ordinarity assertion.
\end{lemma}

\begin{proof}
The pair intercept is
\[
 \frac{rp'-r'p}{p'-p}=v,
\]
and on the lift determined by the line,
\[
 \left|\frac{r-n}{p}-\frac{r'-n}{p'}\right|
 =\frac{|n-v|\,|p-p'|}{pp'}.
\]
The close-pair condition and $pp'\leq4X^2$ give
\eqref{eq:intercept-vacuity}.  Its right-hand side is $2+o(1)$ in the
stated range; integrality gives the assertion.  Each prime has only
$O(1)$ possible partners at those gaps.
\end{proof}

The central Ap\'ery zero, when it occurs, has intercept $-1/2$.
Lemma~\ref{lem:intercept-vacuity} therefore shows that density-one
ordinarity for the non-CM form $8.4.a.a$ is irrelevant to the in-range
problem.  The analogous forced central line for the $\zeta(2)$ Ap\'ery
sequence is likewise an out-of-range obstruction only.

\begin{conjecture}[The first shifted-prime track]
\label{conj:shifted-prime-divisor}
For every fixed $A>0$, uniformly for $y\in\mathcal J_X$,
\[
 \#\{p\in(X,2X]:0\leq y-p<p,\ p\mid b_{y-p}\}
 =o\left(\frac{\pi(X)}{L}\right).
\]
\end{conjecture}

This is the slope-$1$ member of the terminal single-arc family below.
It is a shifted-prime divisor problem for a fixed holonomic sequence,
in the circle of pointwise questions represented by Hooley and by
Murty--Rosen--Silverman \cite{Hooley1973,MurtyRosenSilverman1996}.
It is not used as a substitute for the all-slope collision conjecture.

\begin{proposition}[Dyadic maximum criterion]\label{prop:dyadic-max}
Suppose that, for some fixed $\delta>0$ and every sufficiently large
$X$,
\[
 \max_{m<X^2}K_X(m)\ll X^{1-\delta}.
\]
Then $\log G_n=o(n)$ for every $n$.
More generally, the same conclusion holds if
$|E_X|\ll X^{1-\delta}$ and the displayed maximum bound is known for
\[
 K_X^{G}(m)
 =\#\{p\in(X,2X]\setminus E_X:m\bmod p\in Z_p\}.
\]
\end{proposition}

\begin{proof}
The blocks with $X<2\sqrt n$ cost $O(\sqrt n\log n)$.  On every
remaining block, $n<X^2$, and the assumed estimate gives
$O(X^{1-\delta}\log(2X))$.  The dyadic sum is dominated by the top
block and is $O(n^{1-\delta}\log n)=o(n)$.  The same calculation bounds
the discarded primes without imposing any condition on their residues.
Theorem~\ref{thm:pointwise-reduction} completes the proof.
\end{proof}

\subsection{The depth-$2$ pair gate}

Write
\[
 \mathcal M_X=\sum_{X<p\leq2X}|Z_p|.
\]

\begin{conjecture}[Mass conjecture $(\mathrm C1)$]\label{conj:C1}
For every $\varepsilon>0$,
\[
 \mathcal M_X\ll_{\varepsilon}X^{1+\varepsilon}.
 \tag{C1}\label{eq:C1}
\]
\end{conjecture}

The proved bound in Lemma~\ref{lem:zero-fiber} is
$\mathcal M_X\ll X^{5/3+o(1)}$.  Computations instead suggest about one
root per prime on average: after grouping noncentral roots into
reflection pairs, the pair count is close to a Poisson variable of mean
$1/2$.  Thus $\mathbb E|Z_p|\approx1$, not $1/8$.  The central fixed
point prevents an unconditional parity assertion for $|Z_p|$.

For a set $E_X$ of exceptional block primes define
\[
 P_{2,X}^{G}
 =\sum_{y\in\mathcal J_X}\binom{K_X^{G}(y)}2.
 \tag{6.3}\label{eq:P2-definition}
\]
We write $P_{2,X}$ for the same quantity with $E_X=\varnothing$.

\begin{proposition}[Pair rigidity and the depth-$2$ gate]
\label{prop:pair-gate}
For fixed $A$ and sufficiently large $X$, a pair of distinct block
primes together with one zero class at each prime has at most one
landing in $\mathcal J_X$.  Moreover,
\[
 \sup_{y\in\mathcal J_X}K_X^{G}(y)
 \leq1+\sqrt{2P_{2,X}^{G}}.
 \tag{6.4}\label{eq:pair-gate}
\]
Thus, for any $\delta>0$,
\[
 P_{2,X}^{G}\ll X^{2-2\delta}
 \quad\Longrightarrow\quad
 \sup_{y\in\mathcal J_X}K_X^{G}(y)\ll X^{1-\delta}.
 \tag{6.5}\label{eq:P2-to-sup}
\]
Together with $|E_X|\ll X^{1-\delta}$ at every dyadic scale, this
implies the pointwise conclusion.
\end{proposition}

\begin{proof}
Two simultaneous zero-class congruences have period $pq$.  Since
$p,q>X$, one has $pq>X^2>XL=\max\mathcal J_X$ for sufficiently large
$X$, so two solutions cannot lie in the window.  For each fixed $y$,
$\binom{K_X^{G}(y)}2\leq P_{2,X}^{G}$, which gives
\eqref{eq:pair-gate}.  The last assertion is
Proposition~\ref{prop:dyadic-max}.
\end{proof}

\begin{conjecture}[Exceptional-tolerant pair gate $(P_2)^G$]
\label{conj:P2G}
For every fixed $A>0$ there are $\delta>0$ and sets $E_X$, chosen
independently of $y$, such that
\[
 |E_X|\ll X^{1-\delta},\qquad
 P_{2,X}^{G}\ll X^{2-2\delta}.
 \tag{$P_2^G$}\label{eq:P2G}
\]
\end{conjecture}

Under $(\mathrm C1)$, Markov's inequality permits the canonical choice
$E_X=\{p:|Z_p|>X^\varepsilon\}$, for which
$|E_X|\leq X^{1-\varepsilon+o(1)}$.  Also under $(\mathrm C1)$, the
existence of some power saving in \eqref{eq:P2G} is equivalent, after
renaming the exponent, to a power-saving bound for
$\sup K_X^G$.  Indeed,
\[
 \sum_{y\in\mathcal J_X}K_X^G(y)\leq L\mathcal M_X,
 \qquad
 P_{2,X}^{G}\leq\frac12(\sup K_X^G)
       \sum_yK_X^G(y).
\]
This is the exact depth-$2$ target.  It requires neither a third moment
nor an input of Elliott--Halberstam strength.

\subsection{Single arcs and exact collisions}

For $1\leq a\leq L$ and $y\in\mathcal J_X$, define
\begin{align*}
 S_{a,X}(y)
 &=\{p\in(X,2X]:0\leq y-ap<p,\ y-ap\in Z_p\},\\
 R_{a,X}(y)&=|S_{a,X}(y)|.
\end{align*}
The digit restriction is part of the definition.  It gives
\[
 K_X(y)=\sum_{1\leq a\leq L}R_{a,X}(y),
 \qquad
 \sum_{y\in\mathcal J_X}R_{a,X}(y)\leq\mathcal M_X.
 \tag{6.6}\label{eq:arc-mass}
\]
Cauchy's inequality and \eqref{eq:P2-definition} yield
\[
 P_{2,X}
 \leq\frac12\sum_yK_X(y)^2
 \leq\frac{L^2}{2}\max_{a\leq L}
       \sum_{y\in\mathcal J_X}R_{a,X}(y)^2.
 \tag{6.7}\label{eq:single-arc-collapse}
\]
The Gessel--Coster Lucas congruence gives, for $p\in S_{a,X}(y)$,
\[
 b_y=b_{ap+(y-ap)}\equiv b_a b_{y-ap}\equiv0\pmod p.
\]
Consequently
\[
 R_{a,X}(y)\leq
 \omega_{(X,2X]}(b_y),
 \tag{6.8}\label{eq:arc-omega}
\]
where $\omega_I(m)$ counts distinct prime divisors of $m$ in $I$.

\begin{conjecture}[Single-arc cap $(\mathrm{SA})$]
\label{conj:SA}
For every fixed $A>0$ there is $\eta>0$ such that, uniformly for
$a\leq L$ and $y\in\mathcal J_X$,
\[
 R_{a,X}(y)\ll_A X^{1-\eta}.
 \tag{$\mathrm{SA}_\eta$}\label{eq:SA}
\]
\end{conjecture}

\begin{conjecture}[Middle-prime cap $(\Omega)$]
\label{conj:Omega}
For every fixed $A>0$ there is $\eta>0$ such that, uniformly for
$y\in\mathcal J_X$,
\[
 \omega_{(X,2X]}(b_y)\ll_A X^{1-\eta}.
 \tag{$\Omega_\eta$}\label{eq:Omega}
\]
\end{conjecture}

In the existential ``some positive power'' sense,
Conjectures~\ref{conj:SA} and~\ref{conj:Omega} are equivalent after
renaming $\eta$.  One direction is \eqref{eq:arc-omega}.  Conversely,
Lucas splits a prime counted by $\omega_{(X,2X]}(b_y)$ between the
lower digit and the condition $p\mid b_{\lfloor y/p\rfloor}$.  The
latter channel is $O(y^{2/3})$ by Lemma~\ref{lem:digit-channels}, while
the lower channel is at most $L\max_aR_{a,X}(y)$.

Define the ordered exact-collision count
\begin{align}
 N_{\mathrm{coll},a}(X)
 &=\#\bigl\{(p,z;q,w):p\ne q,\ p,q\in(X,2X],
       \ z\in Z_p,\ w\in Z_q,\notag\\[-2pt]
 &\hspace{92pt}ap+z=aq+w\in\mathcal J_X\bigr\}.
 \tag{6.9}\label{eq:Ncoll-definition}
\end{align}
Opening the square gives the exact identity
\[
 \sum_{y\in\mathcal J_X}R_{a,X}(y)^2
 =\sum_{y\in\mathcal J_X}R_{a,X}(y)
  +N_{\mathrm{coll},a}(X).
 \tag{6.10}\label{eq:Ncoll-identity}
\]

\begin{conjecture}[Terminal exact-collision conjecture]
\label{conj:Ncoll}
For every fixed $A>0$ there is $\eta>0$ such that
\[
 \max_{1\leq a\leq L}N_{\mathrm{coll},a}(X)
 \ll_A X^{2-\eta}.
 \tag{$\mathrm{NC}_\eta$}\label{eq:Ncoll}
\]
\end{conjecture}

\begin{theorem}[Exceptional-tolerant reduction to exact collisions]
\label{thm:exact-collision-reduction}
Conjecture~\ref{conj:Ncoll} implies the pointwise conclusion
$\log G_n=o(n)$.  Under Conjecture~\ref{conj:C1}, the three properties
$(\mathrm{SA})$, $(\Omega)$, and $(\mathrm{NC})$ are equivalent in the
sense that each asserts the existence of some fixed positive power
saving; their exponents need not be the same.  They are also equivalent,
in this sense, to the power-saving target for $\sup K_X$ and to
$(P_2)^G$.
\end{theorem}

\begin{proof}
From \eqref{eq:Ncoll-definition}, for each $a,y$,
\[
 R_{a,X}(y)(R_{a,X}(y)-1)
 \leq N_{\mathrm{coll},a}(X).
\]
Thus $(\mathrm{NC}_\eta)$ gives
$R_{a,X}(y)\ll X^{1-\eta/2}$.  Conversely, $(\mathrm{SA}_\eta)$ and
$(\mathrm C1)$ give
\[
 N_{\mathrm{coll},a}(X)
 \leq(\sup_yR_{a,X}(y))\sum_yR_{a,X}(y)
 \ll X^{2-\eta+o(1)}.
\]
This proves the asserted equivalence, together with the Lucas-channel
comparison above.  Finally, \eqref{eq:single-arc-collapse},
\eqref{eq:Ncoll-identity}, and the proved estimate
$\mathcal M_X\ll X^{5/3+o(1)}$ give a positive power saving for
$P_{2,X}$.  Proposition~\ref{prop:pair-gate} and then
Proposition~\ref{prop:dyadic-max} prove the pointwise conclusion.
\end{proof}

\begin{remark}[No exceptional-$y$ loophole]
An ``off a sparse set of $y$'' formulation is insufficient for a
pointwise theorem unless the exceptional fibers receive their own
collision bound: one exceptional $y$ may carry a fiber of size
$X^{1-o(1)}$.  The uniform formulation \eqref{eq:SA}, or equivalently
the global moment formulation \eqref{eq:Ncoll}, is the precise terminal
statement used here.
\end{remark}

\begin{remark}[Numerical and structural evidence for the terminal conjecture]
\label{rem:evidence}
Direct computation supports \eqref{eq:Ncoll}.  With $G_n$ computed
exactly for all $n\leq 15000$, the extreme values obey the
Poisson-independence law
\[
 \max_{n\leq N}\log G_n\;\approx\;0.85\,(\log N)^2,
\]
a single constant across three orders of magnitude, so
$\log G_n/n=O((\log n)^2/n)$ with no phase transition.  The two channels
of a hypothetical counterexample are both closed numerically.  The
small-prime part is unconditionally $O(\sqrt n\log n)$, so only the
middle collision count can grow; but the exact avalanche count
$K(n)=\#\{\sqrt n<p\leq n:p\mid G_n\}$ never exceeds $7$ for $n\leq15000$
(a counterexample needs $K(n)\gg n/\log n$), and the middle-prime
multiplicity is identically $1$ outside the range $n\leq 6$, so the
full-multiplicity form \eqref{eq:T} shows no phantom.  A targeted search
of the one structured locus is equally negative: the non-ordinary primes
$p\mid\gamma(p)$ of $8.4.\mathrm{a}.\mathrm{a}$ below $12010$ are only
$\{11,3137\}$, and at their forced apex indices $d=(p-1)/2$ the ratio
$\log G_d/d$ is unremarkable while the forced prime does not divide
$G_d$; the coincidence of Lemma-\ref{lem:intercept-vacuity} type lives on
an intercept-$O(1)$ line and does not concentrate at a single $n$.  This
is empirical and heuristic support, not a proof of \eqref{eq:Ncoll}.
\end{remark}

\section{Unconditional structural theorems}\label{sec:structure}

\subsection{The gap eliminant $U_g$}\label{sec:gap-eliminant}

Fix integers $g\geq1$ and $r\geq0$.  Let $u^{(r)}$ be the rational
solution of the Ap\'ery recurrence initialized by
\[
 u^{(r)}_r=0,
 \qquad
 u^{(r)}_{r+1}=1.
\]
For concision, write $u_j=u^{(r)}_j$ while $r$ is fixed.

\begin{lemma}[Eliminant identity]\label{lem:eliminant-identity}
Let $p\geq7$ and $0\leq r<r+g<p$.  If $r\in Z_p$, then
$p\nmid b_{r+1}$ and
\[
 p\mid b_{r+g}
 \quad\Longleftrightarrow\quad
 p\mid u_{r+g}.
\]
Here the last divisibility means that the $p$-integral rational number
$u_{r+g}$ reduces to zero in $\mathbb F_p$.
\end{lemma}

\begin{proof}
Lemma~\ref{lem:no-consecutive} gives $p\nmid b_{r+1}$.  On the interval
$[r,r+g]$ all denominators in the recurrence are $p$-units.  The two
solutions $b_j$ and $b_{r+1}u_j$ have the same initial values at
$j=r,r+1$, so uniqueness for the recurrence gives
\[
 b_{r+j}\equiv b_{r+1}u_{r+j}\pmod p
 \qquad(0\leq j\leq g).
\]
The assertion follows because $b_{r+1}$ is a unit.
\end{proof}

Recall the continuants $N_g(x)$ and $Q_g(x)$ from
\eqref{eq:solution-formula}.  That formula, applied to the normalized
solution above, gives
\[
 Q_g(r)u_{r+g}=N_g(r),
 \qquad
 Q_g(r)=\prod_{j=2}^g(r+j)^3.
\]

\begin{theorem}[Gap-eliminant theorem]\label{thm:Ug}
For every $g\geq1$ there is a fixed polynomial
\[
 U_g(x):=N_g(x)\in\Z[x]
\]
with the following properties:
\begin{enumerate}
\item $\deg U_g=3(g-1)$;
\item the coefficient height of $U_g$ is
$\exp(O(g\log g))$;
\item $U_g(r)>0$ for every integer $r\geq0$;
\item if $p\geq7$ and $0\leq r<r+g<p$, then
\[
 r,r+g\in Z_p\quad\Longrightarrow\quad p\mid U_g(r).
\]
\end{enumerate}
For $g=1$, this means $U_1=1$ and the last assertion is vacuous by
Lemma~\ref{lem:no-consecutive}.
\end{theorem}

\begin{proof}
The identity preceding the theorem proves that clearing the denominators
of $u_{r+g}$ gives the fixed integer polynomial $U_g=N_g$.

Let $L_g$ be the leading coefficient of $N_g$.  From the continuant
recurrence,
\[
 L_1=1,\qquad L_2=34,\qquad L_{g+1}=34L_g-L_{g-1}.
\]
Inductively $L_g>L_{g-1}>0$.  Hence no top-degree cancellation occurs,
and $\deg N_g=3(g-1)$.

For the height, let $\|F\|_1$ denote the sum of the absolute values of
the coefficients of $F$.  Shifting the fixed cubic $P$ and the sixth
power in the continuant recurrence gives
\[
 \|N_{g+1}\|_1
 \leq C(g+1)^6\max(\|N_g\|_1,\|N_{g-1}\|_1).
\]
Iteration yields
$\log\|N_g\|_1=O(\sum_{j\leq g}\log j)=O(g\log g)$, which proves the
height bound.

For positivity, the Wronskian identity gives
\[
 u_{r+g}
 =\frac{a_{r+g}b_r-b_{r+g}a_r}
 {a_{r+1}b_r-b_{r+1}a_r}.
\]
The denominator is $W_r>0$, and the numerator is positive because the
positive sequence $b_j$ has strictly increasing ratios $a_j/b_j$, again
by $W_j>0$.  Thus $u_{r+g}>0$ and consequently
$U_g(r)=Q_g(r)u_{r+g}>0$.

Finally, if $r,r+g\in Z_p$, Lemma~\ref{lem:eliminant-identity} gives
$u_{r+g}\equiv0\pmod p$.  Every factor of $Q_g(r)$ is a $p$-unit, so
$p\mid U_g(r)$.
\end{proof}

\begin{corollary}[Bounded-gap pairs]\label{cor:bounded-gap}
Uniformly for $1\leq G\leq X$,
\[
 \#\left\{(p,r,g):
 \begin{array}{l}
 X<p\leq2X,\quad1\leq g\leq G,\\
 0\leq r<r+g<p,\quad r,r+g\in Z_p
 \end{array}\right\}
 \ll G^2X.
\]
\end{corollary}

\begin{proof}
For $r<2X$ and $g\leq G\leq X$, the degree and height bounds give
\[
 \log U_g(r)\ll g\log X.
\]
Since $U_g(r)>0$, it has $O(g)$ prime divisors exceeding $X$.  Sum this
bound over $r<2X$ and $g\leq G$.
\end{proof}

\begin{remark}[Scope of root distribution]
Theorem~\ref{thm:Ug} embeds bounded-gap double zeros into roots of a
fixed polynomial, which is the algebraic prerequisite for the
Hooley--DFI circle.  It does not by itself prove that the Ap\'ery subset
of those roots is equidistributed.  In particular,
\[
 U_2(r)=(2r+3)(17r^2+51r+39)
\]
has a linear factor.  The theorem of Duke, Friedlander, and Iwaniec
treats roots of an irreducible quadratic to prime moduli \cite{DFI1995},
while Hooley's general theorem orders roots by all moduli
\cite{Hooley1964}.  Thus the
stronger prime-axis equidistribution claim in the working notes is not an
unconditional corollary and is not used here.
\end{remark}

\subsection{A cross-arc eliminant}

\begin{lemma}[Cross-arc lemma]\label{lem:cross-arc}
Fix $A>0$.  Uniformly for $1\leq a,m\leq L$ and
$y,y+m\in\mathcal J_X$, one has, for sufficiently large $X$,
\[
 |S_{a,X}(y)\cap S_{a,X}(y+m)|
 \leq\omega_{>X}(U_m(y))\leq4m.
 \tag{7.1}\label{eq:cross-arc}
\]
\end{lemma}

\begin{proof}
If $p$ belongs to the intersection, then
$r=y-ap$ and $r+m=y+m-ap$ both lie in $Z_p$, with
$0\leq r<r+m<p$.  Theorem~\ref{thm:Ug} gives
\[
 p\mid U_m(r),\qquad U_m(r)\equiv U_m(y)\pmod p.
\]
Thus every such $p$ divides the fixed positive integer $U_m(y)$.
The degree and coefficient-height estimates in
Theorem~\ref{thm:Ug}, together with $y\leq XL$ and $m\leq L$, give
\[
 \log U_m(y)\leq(3+o(1))m\log X.
\]
It therefore has at most $4m$ prime divisors exceeding $X$ for all
sufficiently large $X$.
\end{proof}

\subsection{Structured resonance}

For $R,H<X/2$, let $Z_p^{\rm str}(R,H)$ consist of those $r\in Z_p$
which either satisfy
\[
 \min(r,p-1-r)\leq R
 \quad\text{or}\quad
 |2r-(p-1)|\leq2R,
 \tag{7.2}\label{eq:anchors}
\]
or have a second zero $s\in Z_p$ with $0<|r-s|\leq H$.  Put
\[
 K_{R,H}(x)=\#\{p\in(X,2X]:x\bmod p\in Z_p^{\rm str}(R,H)\}.
\]

\begin{theorem}[Structured-resonance theorem]
\label{thm:structured-resonance}
Let $0\leq\delta_0<1$, $2X\leq x\leq Y\leq X^{2+\delta_0}$, and
$1\leq R,H\leq X/4$.  Then
\[
 K_{R,H}(x)\ll_{\delta_0}R+H^2+1.
 \tag{7.3}\label{eq:structured-K}
\]
Consequently, with $(t)_3=t(t-1)(t-2)$,
\[
 A_3^{\rm str}(Y,X;R,H)
 :=\sum_{2X\leq x\leq Y}(K_{R,H}(x))_3
 \ll_{\delta_0}Y(R+H^2+1)^3.
 \tag{7.4}\label{eq:structured-A3}
\]
The contribution in which all three primes hit the exact center
$r=(p-1)/2$ is zero for sufficiently large $X$.
\end{theorem}

\begin{proof}
The anchor conditions force a block prime to divide the nonzero integer
\[
 L_R(x)=
 \prod_{0\leq t\leq R}(x-t)(x+1+t)
 \prod_{|t|\leq R}(2x+1-2t).
\]
There are $O(R+1)$ linear factors, and
$\log|L_R(x)|\ll_{\delta_0}(R+1)\log X$; hence only
$O_{\delta_0}(R+1)$ prime divisors of $L_R(x)$ exceed $X$.

If $x\bmod p$ has a second zero at distance $h\leq H$, then, according
to which zero is smaller, Theorem~\ref{thm:Ug} makes $p$ divide either
$U_h(x)$ or $U_h(x-h)$.  Thus every close-pair prime divides
\[
 C_H(x)=\prod_{1\leq h\leq H}U_h(x)U_h(x-h),
\]
which is positive because $x-h>0$.  The degree and height of the
factors give
\[
 \log C_H(x)\ll_{\delta_0}H^2\log X.
\]
It has $O_{\delta_0}(H^2)$ prime divisors exceeding $X$.  This proves
\eqref{eq:structured-K}; summing its cube proves
\eqref{eq:structured-A3}.

Finally, an exact-center hit gives $p\mid2x+1$.  Three distinct block
primes would give $pqs\mid2x+1$, whereas
$pqs>X^3>2X^{2+\delta_0}+1\geq2x+1$ for sufficiently large $X$.
This is impossible.
\end{proof}

\subsection{Two lenses on the zero set}

\begin{remark}[The Ap\'ery--Hasse polynomial]
For a prime $p$, define
\[
 \Phi_p(T)=\sum_{0\leq k<p}
   \binom Tk^2\binom{T+k}k^2\in\mathbb F_p[T].
 \tag{7.5}\label{eq:Hasse-polynomial}
\]
Then $\deg\Phi_p=4(p-1)$ and
\[
 Z_p=\{r\in\{0,\ldots,p-1\}:\Phi_p(r)=0\}.
\]
The classical identity
\[
 b_r={}_4F_3\!\left[
 \begin{matrix}-r,-r,r+1,r+1\\1,1,1\end{matrix};1\right]
\]
has finite-field and $p$-adic counterparts.  In particular, for
$1\leq r\leq p-2$, Kalita--Chetry express $b_r$ modulo $p$ as a
Greene ${}_4F_3$ character sum, while McCarthy's notation writes the
same balanced family as
${}_4G(r/(p-1),r/(p-1),1-r/(p-1),1-r/(p-1))_p$
\cite{KalitaChetry2017,Greene1987,McCarthy2012}.  These are sums over
characters, not bounded products, so they supply a moment lens rather
than a compressive cross-prime invariant.

At the central index $r=(p-1)/2$, Ahlgren--Ono and Kilbourn give the
modular-form congruences with the weight-$4$, level-$8$ newform
$\eta(2\tau)^4\eta(4\tau)^4$ \cite{AhlgrenOno2000,Kilbourn2006}.
The Ap\'ery index is $(p-1)/2$, not $p-1$; indeed
$b_{p-1}\equiv1\pmod p$.
\end{remark}

\begin{remark}[Projective Frobenius lens]
Let $m=(p-1)/2$ and let $\gamma(p)$ be the $p$th coefficient of
$\eta(2\tau)^4\eta(4\tau)^4$.  Reducing the recurrence around $m$
produces universal rational continuant coefficients $A_j,B_j$, on
every local interval where the recurrence denominators are $p$-units,
such that
\[
 \begin{aligned}
 b_{m+j}&\equiv A_j\gamma(p)+B_jb_{m-1}\pmod p,\\
 (A_0,B_0)&=(1,0),\qquad(A_1,B_1)=(0,1).
 \end{aligned}
 \tag{7.6}\label{eq:two-seed}
\]
The Wronskian prevents the two seeds from vanishing simultaneously.
Thus the root condition can be read projectively as
\[
 [A_j:B_j]=u_p,
 \qquad
 u_p=[-b_{m-1}:\gamma(p)]\in\mathbb P^1(\mathbb F_p).
 \tag{7.7}\label{eq:projective-frobenius}
\]
This is complementary to \eqref{eq:Hasse-polynomial}: a fixed
continuant curve meets a moving projective Frobenius datum.  No
cross-prime anti-concentration theorem for $u_p$ is asserted here.
Newton--Thorne symmetric-power functoriality controls the archimedean
Sato--Tate distribution, not the self-characteristic condition
$p\mid\gamma(p)$ \cite{NewtonThorne2021}.
\end{remark}

\section{Difficulty certification}\label{sec:difficulty}

\subsection{Phantom locality}

\begin{remark}[Phantom-locality doctrine]\label{rem:phantom-locality}
Every currently available holonomic divisibility constraint in this
problem is local at one prime: the Lucas and Dwork congruences,
reflection, the Wronskian and Casoratian, the $U_g$ and cross-arc
certificates, coefficient heights, and the polynomial $\Phi_p$.
None supplies a map from zero positions at different primes to a common
object with substantially fewer fibers.  They therefore do not exclude
the possibility that, at many primes, one locally admissible zero lies
at $z_p\equiv y_0-ap\pmod p$.  Reflection adds a partner locally and the
eliminants constrain local pairs, but neither operation couples the
chosen positions across different primes.  Such an abstract arc
phantom would create one large exact fiber.

This is a limitation of the banked invariant list, not a construction
of alternative roots of the actual polynomial $\Phi_p$.  In
particular, it does not prove logical independence from all possible
Ap\'ery identities.  It identifies the missing kind of input: a
cross-prime compressive invariant.  The multiplicative-order fibration
in Lemma~\ref{lem:Mersenne} below is the model of such an invariant.
\end{remark}

The usual single-value eliminant is $b_y$ itself, whose exponential
height reproduces the trivial divisor bound.  Hermite--Pad\'e
denominator bounds, the Casoratian, and fixed-prime hypergeometric
identities do not alter this cross-prime obstruction.  The available
Ap\'ery factorization and congruence results
\cite{CFVMZ2026,Delaygue2018,RowlandYassawiKrattenthaler2021}
provide strong local information but no such fibration.

\subsection{A fibration calibration}

\begin{lemma}[Mersenne fibration]\label{lem:Mersenne}
Fix $A>0$.  Uniformly for $1\leq a\leq L$ and
$y\in\mathcal J_X$, the admissible Mersenne arc satisfies
\[
 \#\{p\in(X,2X]:0\leq y-ap<p,\ p\mid2^{y-ap}-1\}
 \ll_A X^{1/2+o(1)}.
 \tag{8.1}\label{eq:Mersenne}
\]
\end{lemma}

\begin{proof}
Let $d_p=\operatorname{ord}_p(2)$.  The divisibility condition is
$d_p\mid y-ap$.  Since $d_p\mid p-1$, it follows that
\[
 d_p\mid y-a.
 \tag{8.2}\label{eq:Mersenne-fibration}
\]
Here $0<y-a\ll_A X(\log X)^A$, so
$\tau(y-a)=X^{o(1)}$ uniformly.

For a divisor $d\leq\sqrt X$ of $y-a$, every prime with $d_p=d$
divides $2^d-1$.  The number of its prime divisors exceeding $X$ is
$O(d/\log X)$.  Summing over the possible $d$ costs
$X^{1/2+o(1)}$.  If $d>\sqrt X$, then $p\equiv1\pmod d$, so the
interval $(X,2X]$ contains $O(X/d+1)=O(\sqrt X)$ candidates for each
such divisor.  Another use of $\tau(y-a)=X^{o(1)}$ proves
\eqref{eq:Mersenne}.
\end{proof}

The restriction to the in-range window is essential: with unrestricted
$y$, the integer $y-a$ could be a common multiple of all relevant
orders.  The lemma exhibits exactly what Ap\'ery lacks.  Mersenne zero
positions fiber through the divisor lattice of one small integer
$y-a$; Ap\'ery zero positions are not known to fiber through a subgroup
or any other compressive cross-prime object.

\subsection{The horizontal Fourier barrier}

Put
\[
 W_X(h)=\sum_{X<p\leq2X}\ \sum_{m\in Z_p}e\left(\frac{hm}{p}\right),
 \qquad
 \mathcal S_X=\{p\in(X,2X]:p\mid\gamma(p)\},
 \tag{8.3}\label{eq:W-definition}
\]
where $\gamma(p)$ is the coefficient in
$\eta(2\tau)^4\eta(4\tau)^4$.  Ahlgren--Ono identify
\[
 p\in\mathcal S_X
 \quad\Longleftrightarrow\quad
 (p-1)/2\in Z_p.
 \tag{8.4}\label{eq:central-equivalence}
\]

\begin{theorem}[Barrier R9.2: Lang--Trotter hardness]
\label{thm:R92}
Let $2\leq H\leq X$.  There is an absolute $c>0$ such that
\[
 cH|\mathcal S_X|
 \leq\mathcal M_X+2\sum_{1\leq h<H}|W_X(h)|.
 \tag{8.5}\label{eq:Fejer-barrier}
\]
Consequently, if for some fixed $\delta>0$ one has, with
$H=\lfloor X/(4L)\rfloor$,
\[
 \max_{1\leq h<H}|W_X(h)|\ll X^{1-\delta},
 \tag{8.6}\label{eq:W-sup-bound}
\]
then $|\mathcal S_X|\ll X^{1-\delta'}$ for some $\delta'>0$.
The same conclusion follows from
\[
 \sum_{1\leq h<H}|W_X(h)|^2\ll H X^{2-2\delta}.
 \tag{8.7}\label{eq:W-second-bound}
\]
In particular, either horizontal estimate implies that the primes
$p$ for which $p\mid\gamma(p)$ have density zero.
\end{theorem}

\begin{proof}
For the Fej\'er kernel
\[
 F_H(t)=\frac1H\left|\sum_{j=0}^{H-1}e(jt)\right|^2
 =\sum_{|h|<H}\left(1-\frac{|h|}{H}\right)e(ht),
\]
positivity gives
\begin{align*}
 0&\leq
 \sum_{|h|<H}\left(1-\frac{|h|}{H}\right)
 e(-h/2)W_X(h)\\
 &=\sum_{X<p\leq2X}\ \sum_{m\in Z_p}
 F_H\left(\frac mp-\frac12\right).
\end{align*}
If $m=(p-1)/2$ and $H\leq X<p$, then
$F_H(-1/(2p))\geq cH$.  Keeping these nonnegative terms and bounding
the nonzero Fourier modes in absolute value proves
\eqref{eq:Fejer-barrier}.

Lemma~\ref{lem:zero-fiber} gives
$\mathcal M_X\ll X^{5/3+o(1)}$.  Divide
\eqref{eq:Fejer-barrier} by $H\asymp X/L$ and use
\eqref{eq:W-sup-bound}; the two terms are
$X^{2/3+o(1)}$ and $X^{1-\delta}$, up to powers of $\log X$.
This gives a fixed power saving.  Under \eqref{eq:W-second-bound}, use
Cauchy's inequality before dividing by $H$.
\end{proof}

The form $8.4.a.a$ is non-CM of weight $4$.  No theorem currently
gives even $o(\pi(X))$ non-ordinary primes for this fixed form.  Deligne's
bound allows $\gamma(p)=cp$ with $|c|\ll\sqrt p$ and does not force
vanishing, while Newton--Thorne controls an archimedean Sato--Tate
distribution rather than divisibility by the varying prime.  Thus
Theorem~\ref{thm:R92} is an implication to a recognized open
Lang--Trotter-type problem, not a claim that \eqref{eq:W-sup-bound} is
logically unprovable.

\begin{remark}[Barrier R9.1: bandwidth insufficiency]
\label{rem:R91}
For a planted arc $m_p\equiv y_0\pmod p$, the summand in $W_X(h)$ is
$e(hy_0/p)$.  Reciprocal-phase cancellation can make the high-frequency
range look equidistributed while all planted primes still collide at
the exact integer $y_0$.  The working-note R9.1 proposes correcting the
low frequencies by a small perturbation among the remaining primes,
thereby producing an abstract full-bandwidth phantom.  The notes do not
supply a positivity- and integrality-preserving realization as actual
root sets of the fixed polynomials $\Phi_p$, so this stronger
construction is recorded as a barrier model rather than a theorem.

Its robust conclusion is categorical: Fourier bounds for the normalized
points $m/p$ control intervals of positive width.  The condition
$ap+m=aq+w$ is an exact equality at scale zero.  Without an input tied
to that equality, horizontal equidistribution is neither the terminal
statement nor a replacement for \eqref{eq:Ncoll}.
\end{remark}

\subsection{The scale-zero core}

Identity~\eqref{eq:Ncoll-identity} is the corrected detector.  Opening
$N_{\mathrm{coll},a}$ by the Chinese remainder theorem and applying
additive orthogonality produces a reciprocity-twisted quadrilinear form
with moduli $pq$.  Away from the locus on which the two reciprocal
phases recombine, this has the formal shape treated by the bilinear and
trilinear Kloosterman-fraction estimates of Duke--Friedlander--Iwaniec
and Bettin--Chandee \cite{DFI1997,BettinChandee2018}.  No estimate for
the Ap\'ery-weighted form is proved here; attaching those tools to the
non-degenerate periphery remains a prospective step.

\begin{remark}[Degeneracy certificate]
On the recombination locus, an arc phantom makes both phases combine
to a constant depending on the common exact value.  The mass on that
locus is precisely $N_{\mathrm{coll},a}$, so an equidistribution
estimate for the complementary locus leaves the terminal conjecture
unchanged.  Rationally aligned central clusters do not import the
weight-$4$ ordinarity obstruction into $N_{\mathrm{coll},a}$: after
clearing the fixed denominator, the collision congruence is a multiple
of $pq$ whose absolute value is smaller than $pq$, and hence is an
integer equality.  This is why the horizontal Fourier route asks for
strictly more than the exact-collision reduction needs.
\end{remark}

\begin{remark}[Current technology boundary]
The nearest exact-divisibility tools exploit multiplicative structure:
large prime factors in Lucas and Lehmer sequences
\cite{Stewart2012}, or Subspace-Theorem bounds for greatest common
divisors of $S$-unit values \cite{CorvajaZannier2005}.  The available
results on Ap\'ery numbers are one-sided valuation or Lucas-congruence
statements \cite{LucaShparlinski2008,Delaygue2018,
RowlandYassawiKrattenthaler2021}; they do not give an upper bound for
$\omega_{(X,2X]}(b_y)$.  No standard GRH-type hypothesis presently
formulates, let alone implies, \eqref{eq:Omega}.

The two live directions are therefore: construct a cross-prime
compressive invariant for the vanishing positions of the mod-$p$
Picard--Fuchs solutions, or develop an exact-collision tool for
holonomic values.  The non-degenerate Kloosterman-fraction periphery
\cite{DFI1997,BettinChandee2018} is the natural first
test case; matching the Ap\'ery weights to the hypotheses of those
methods, so as to obtain a theorem-level bound on the provable
periphery of $(\Omega)$, remains to be carried out.
\end{remark}

\subsection{The six-fold barrier map}\label{sec:sixfold-barrier}

The preceding barriers admit a common organization.  The word
``complete'' below is relative to the audited mechanisms: five
ensemble methods and one single-index adelic method.  A barrier
certificate says that the stated input class cannot distinguish the
terminal event; it is not a logical independence result for arbitrary
future mathematics.

\subsubsection{Sieve complexity and the integrality criterion}

For a family of $p$-local relations $\mathcal R_p$, define its
\emph{sieve complexity} $c_X(\mathcal R)$ to be the least number of
global integral relations whose reductions cover $\mathcal R_p$ for
all $p\in(X,2X]$.  Only the order of magnitude of this invariant will
be used.  The calibrating examples are
\[
\begin{array}{c|c|c}
\text{relation}&c_X&\text{status of the corresponding cover}\\ \hline
\text{weight-$3$ ordinarity}&O(1)&\text{provable}\\
\text{weight-$4$ form $8.4.a.a$}&X^{1/2}&\text{open}\\
\text{Mersenne divisibility}&X^{o(1)}&\text{provable}\\
\text{Ap\'ery divisibility at index }n&e^{\Theta(n)}&
  \text{outside the subexponential range}.
\end{array}
\tag{8.8}\label{eq:sieve-complexity-table}
\]
For the first row, Deligne's interval $|a_p|\leq2p$ leaves the five
integer classes $a_p=cp$.  For weight $4$, the interval
$|\gamma(p)|\leq2p^{3/2}$ leaves $\asymp\sqrt X$ such classes.  In the
Mersenne model, the classes are the divisors of $y-a$, as in
\eqref{eq:Mersenne-fibration}.  For Ap\'ery numbers, writing a
divisibility relation as a global equality introduces a cofactor of
exponential range.

\begin{proposition}[Integrality criterion]
\label{prop:integrality-criterion}
A fixed-quotient Galois or sieve argument can test a relation uniformly
across varying primes only after that relation has been expressed by
global integral classes.  A condition defined only in $\mathbb F_p$
has no reduction in a fixed auxiliary quotient.  Consequently, such a
method applied to $n\bmod p\in Z_p$ must first pay the covering cost
$c_X$ in \eqref{eq:sieve-complexity-table}.
\end{proposition}

\begin{proof}[Proof sketch]
A global equality over $\mathbb Z$ may be reduced modulo every fixed
auxiliary prime and is therefore visible to a fixed-quotient sieve.
By contrast, the residue $z_p\in\mathbb F_p$ belongs to a different
fiber for each $p$; there is no canonical image of the equality
$n\equiv z_p\pmod p$ in a fixed quotient.  Replacing it by global
equalities amounts precisely to enumerating its integral lifts.  The
four counts following the table give the asserted calibration.  Thus
the criterion does not say that large image never detects equality;
it separates globally integral equalities, which it can detect, from
self-characteristic $p$-local equalities, which require a cover.
\end{proof}

\subsubsection{The five ensemble barriers}

\begin{proposition}[Barrier R9: analytic closure]
\label{prop:barrier-analytic}
The horizontal Fourier route is closed in two complementary senses.
First, the estimates \eqref{eq:W-sup-bound} or
\eqref{eq:W-second-bound} imply a density-zero theorem for the
weight-$4$ nonordinary primes through Theorem~\ref{thm:R92}, an open
Lang--Trotter-type endpoint.  Second, the bandwidth model in
Remark~\ref{rem:R91} can preserve an exact planted collision while
matching the available positive-width distribution data.  In
particular, the exact-collision count \eqref{eq:Ncoll} is a scale-zero
statistic not determined by horizontal equidistribution.
\end{proposition}

\begin{proof}[Proof sketch]
The first assertion is the Fej\'er-positivity argument
\eqref{eq:Fejer-barrier}.  The second is the reciprocal-phase
calculation in Remark~\ref{rem:R91}.  That calculation is a barrier
model, not a realization by the actual roots of $\Phi_p$; the robust
conclusion is only that interval discrepancy and exact equality belong
to different categories of information.
\end{proof}

\begin{proposition}[Barrier R10a: collision decoupling]
\label{prop:collision-decoupling}
For distinct $p,q\in(X,2X]$, an exact collision factors into
\[
 \bigl[z\in Z_p\bigr]\quad\times\quad
 \bigl[w\in Z_q\bigr]\quad\times\quad
 \bigl[ap+z=aq+w\in\mathcal J_X\bigr].
 \tag{8.9}\label{eq:collision-decoupling}
\]
The first two factors are separate self-characteristic conditions,
while the last factor is a Frobenius-free CRT window condition.  There
is no joint algebraic relation between the projective data $u_p$ and
$u_q$ in \eqref{eq:collision-decoupling}.  Accordingly, joint
large-image, Mumford--Tate, and fixed-quotient Chebotarev information
does not by itself make the collision rare.
\end{proposition}

\begin{proof}[Proof sketch]
For fixed $p$ and $q$, CRT makes the two residue conditions independent;
the only coupling is the requirement that their common lift lie in the
short archimedean window.  Each local factor lives in its own
characteristic and $b_{(p-3)/2}$ is not a common Weil-number datum.
Even a square-root Chebotarev error for a pair has size
$(pq)^{1/2}\asymp X$; summing such errors over prime pairs is larger
than the target.  Finally, the abstract coherent zero-section from
Remark~\ref{rem:phantom-locality} is compatible with maximal
fixed-quotient and Sato--Tate data while failing $(\mathrm{SA})$.
This is the Galois-phantom certificate.  As before, it is a certificate
for the stated hypothesis class, not a second system of roots for
$\Phi_p$.
\end{proof}

Put
\[
 \lambda=\log(17+12\sqrt2)=4\log(1+\sqrt2).
\]

\begin{proposition}[Barrier R10b: the $3n$ determinant cap]
\label{prop:determinant-cap}
If an archimedean one-form argument uses an irrationality exponent
$\mu\geq2$ for $\zeta(3)$, then its resulting coefficient is bounded
by the template
\[
 \log G_n\leq
 n\frac{(\mu-1)(3+\lambda)-(\lambda-3)}{\mu}+o(n).
 \tag{8.10}\label{eq:determinant-cap}
\]
At the Rhin--Viola value $\mu=5.51389\ldots$ this is
$5.24673\ldots n$, while the best conceivable limit $\mu\to2$ is
exactly $3n\sim\log d_n$.  Thus neither this one-form method nor its
Hermite--Pad\'e determinant localization can yield $o(n)$ without new
$p$-local input.
\end{proposition}

\begin{proof}[Proof sketch]
The right-hand side of \eqref{eq:determinant-cap} equals
$3+\lambda-2\lambda/\mu$, which is the corrected Rhin--Viola
calculation recorded above.  Sending $\mu$ to the Roth threshold $2$
leaves $3$.  In a multi-form determinant the Ap\'ery row still carries
the denominator $d_n$; the Casoratian has only polynomial height and
prime support below $n$.  Hence localization reproduces, rather than
removes, the same $3n$ budget.  This is the point at which general
Subspace-Theorem and multiplicative-group technology
\cite{Schmidt1980,EvertseSchlickeweiSchmidt2002,
BugeaudCorvajaZannier2003} returns the problem to the missing local
support law.
\end{proof}

Consider the projective skew product over the shift $t\mapsto t+1$
with transfer matrix
\[
 M(t)=
 \begin{pmatrix}
 P(t)&-t^3\\ (t+1)^3&0
 \end{pmatrix}.
 \tag{8.11}\label{eq:mobius-cocycle}
\]
Scalar normalization is irrelevant projectively.

\begin{proposition}[Barrier R11: torus versus $\mathrm{SL}_2$]
\label{prop:barrier-dynamical}
Let $G$ be the difference-Galois group over $\overline{\mathbb Q}(t)$
of the rank-$2$ Ap\'ery difference module defined by
\eqref{eq:mobius-cocycle}.
\begin{enumerate}
\item[(NI-1)] One has $G\supseteq\mathrm{SL}_2$.
\item[(NI-2)] For each fixed degree bound $D$, a nonconstant
invariant $F\in\mathbb F_p(x,t)$ of degree at most $D$ for infinitely
many $p$ would lift to characteristic zero.  Hence, outside finitely
many primes, the mod-$p$ cocycle has no uniformly bounded-complexity
first integral beyond the determinant line, represented here by the
Casoratian.
\end{enumerate}
The Mersenne comparison has toric difference-Galois group and a full
character lattice of invariants.  Thus the torus/$\mathrm{SL}_2$
dichotomy identifies large monodromy with the failure of a compressive
first integral.
\end{proposition}

\begin{proof}[Proof sketch]
Birkhoff asymptotics at infinity leave characteristic factors
$\alpha^{\pm t}t^{-3/2}(1+\cdots)$, with
$\alpha=17+12\sqrt2$.  Substitution in the hyperexponential
possibilities excludes a rank-$1$ solution; the half-integral exponent
is the terminal certificate in this calculation.  A dihedral image
would give a hyperexponential solution of the symmetric-square
recurrence, whose three candidate exponential ratios are
$\alpha^2,1,\alpha^{-2}$.  The terminating Hyper computation excludes
them, in the Hendriks--Singer classification
\cite{HendriksSinger1999}.  Irreducibility, primitivity, and infinitude
then give (NI-1).

For (NI-2), the coefficients of an invariant of degree at most $D$
satisfy a finite linear system over $\mathbb Z[t]$.  Its rank can drop
modulo only the finitely many primes dividing the relevant nonzero
minors.  Invariants for infinitely many primes would therefore lift,
contradicting (NI-1).  Algebraic entropy alone is not the relevant
test: the base shift makes the skew product have zero algebraic entropy
while its invariant ring remains minimal
\cite{BellonViallet1999}.  The Lucas--Gessel factorization at
$t\equiv-1,0\pmod p$ is instead an arithmetic singularity-confinement
event, in the sense of almost-good reduction
\cite{KankiMadaTokihiro2013}; it supplies the already-used digit law,
not a cross-prime first integral.
\end{proof}

\begin{proposition}[Barrier R12: star collapse and the three habitats]
\label{prop:barrier-additive}
Let
\[
 \mathcal T_X=\{(a,p,z):1\leq a\leq L, p\in(X,2X],
 z\in Z_p, ap+z\in\mathcal J_X\}
\]
and let $r_X(y)$ be the number of representations $y=ap+z$ from
$\mathcal T_X$.  Then the following certificates hold.
\begin{enumerate}
\item If $\mathcal A_X=\{ap:1\leq a\leq L, p\in(X,2X]\}$ and
$M=|\mathcal A_X|$, then
\[
 E_+(\mathcal A_X)\gg \frac{M^3}{\log X}.
 \tag{8.12}\label{eq:log-density-energy}
\]
Thus the prime-slope set already has nearly maximal additive energy.
\item The collision graph on $\mathcal T_X$ is a disjoint union of
cliques indexed by $y$, and
\[
 N_*(X):=\sum_y r_X(y)(r_X(y)-1),\qquad
 \sum_y r_X(y)=|\mathcal T_X|.
 \tag{8.13}\label{eq:star-collapse}
\]
Under $(\mathrm C1)$, a power saving for $N_*(X)$ is equivalent, after
renaming exponents, to a power saving for $\max_y r_X(y)$, hence to the
power-saving forms of $(\mathrm{SA})$ and $(\Omega)$.
\item The coherent section $z_p\equiv y_0\pmod p$ has the same
fiber-size, prime-slope, and normalized-distribution inputs as the
abstract phantom, but it creates one fiber of large multiplicity.  In
the line model $v=pu+z_p$, writing
$z_p=y_0-b_pp$ with $b_p=\lfloor y_0/p\rfloor$ makes every line meet
the horizontal fiber $v=y_0$ at an allowed integer $u=b_p$.  After the
section-dependent trivialization $c=u-b_p$, this is the projected
pencil with apex $(0,y_0)$.
\end{enumerate}
Consequently, additive combinatorics has only three relevant habitats:
a fixed auxiliary quotient, closed by
Proposition~\ref{prop:integrality-criterion}; the individual fibers
$\mathbb F_p$, where singleton-scale sets have no internal structure;
and global $\mathbb Z$, where \eqref{eq:log-density-energy} and
\eqref{eq:star-collapse} return the argument to the supremum cap.
Within these hypotheses, sum-product and incidence methods are
phantom-transparent.
\end{proposition}

\begin{proof}[Proof sketch]
Distinct products $ap$ are unique for $a\leq L=o(X)$, so
$M\asymp LX/\log X$, while $\mathcal A_X+\mathcal A_X$ lies in an
interval of length $O(LX)$.  Cauchy's inequality gives
$E_+(\mathcal A_X)\geq M^4/|\mathcal A_X+\mathcal A_X|$, proving
\eqref{eq:log-density-energy}.

Two progressions with distinct moduli $p,q$ can meet at most once in
the window: two meetings would make the differences of their slope
parameters multiples of $q$ and $p$, larger than $L$.  The same-prime
representation is unique.  This proves the clique decomposition.  If
$R=\max_y r_X(y)$ and $T=\sum_y r_X(y)$, then
\[
 R(R-1)\leq N_*(X)\leq RT.
\]
Under $(\mathrm C1)$, $T\leq L\mathcal M_X=X^{1+o(1)}$, giving the
claimed equivalence after exponent renaming.  The third assertion is
the direct substitution $y_0=b_pp+z_p$.  It is a certificate for
hypotheses built from fiber statistics, mass, prime distribution, and
equidistribution; it does not assert that the coherent section is the
actual Ap\'ery zero-section.  Bilinear and incidence estimates can
still treat a nondegenerate periphery
\cite{FouvryIwaniec1983,BombieriFriedlanderIwaniec1986,
DFI1997,BettinChandee2018}, but the concentrated fiber is exactly the
quantity in $(\Omega)$.
\end{proof}

\subsubsection{The single-index adelic barrier}

Let $D_n=\lcm(1,\ldots,n)$, so that $d_n=D_n^3$, and put
\[
 \mathcal B_n=D_n^3b_n,\qquad
 \mathcal A_n=D_n^3a_n,\qquad
 \ell_n=b_n\zeta(3)-a_n.
 \tag{8.14}\label{eq:adelic-data}
\]
The standard Ap\'ery remainders $\ell_n$ are positive and satisfy
$|\ell_n|\asymp e^{-\lambda n}n^{-3/2}$.

\begin{lemma}[Casoratian floor]
\label{lem:casoratian-floor}
For every $n\geq0$,
\[
 |\ell_n|\geq\frac{6}{(n+1)^3b_{n+1}}
 =e^{-\lambda n-O(\log n)}.
 \tag{8.15}\label{eq:casoratian-floor}
\]
\end{lemma}

\begin{proof}
The Casoratian identity gives
\[
 b_{n+1}\ell_n-b_n\ell_{n+1}
 =a_{n+1}b_n-a_nb_{n+1}=\frac6{(n+1)^3}.
\]
Both remainders have the same sign, so
$b_{n+1}|\ell_n|\geq6/(n+1)^3$.  The known growth
$b_{n+1}=e^{\lambda n+O(\log n)}$ proves the last equality.
\end{proof}

\begin{proposition}[Barrier R13a: the critical adelic ledger]
\label{prop:critical-ledger}
For $x_n=(\mathcal B_n:\mathcal A_n)\in\mathbb P^1(\mathbb Q)$,
\[
 h(x_n)=(\lambda+3)n+O(\log n).
\]
The best-form bookkeeping from the Ap\'ery lattice has total local
smallness
\[
 (3\lambda+3)n+O(\log n),
\]
whereas the rank-$2$ Roth--Subspace threshold is
$2h(x_n)=2(\lambda+3)n+O(\log n)$.  The surplus is therefore exactly
\[
 (\lambda-3)n+O(\log n),
 \tag{8.16}\label{eq:critical-ledger}
\]
the same margin used in Ap\'ery's irrationality argument.  It is not a
second budget that can also pay for locating the primes dividing
$b_n$.
\end{proposition}

\begin{proof}[Proof sketch]
The height is the sum of the exponential growth $\lambda n$ of the
Ap\'ery solution and the denominator-clearing tax
$3\log D_n=3n+o(n)$.  Tallying the best local form at infinity and at
the finite places gives $(3\lambda+3)n$; subtracting twice the height
gives \eqref{eq:critical-ledger}.  The only form carrying this positive
surplus at infinity has coefficient $\zeta(3)$ and is therefore not an
algebraic-coefficient form to which the Subspace Theorem applies.
Removing it leaves the ledger exactly critical.  Lemma
\ref{lem:casoratian-floor} also shows that the determinant and
archimedean channels are two readings of one identity and cannot be
added.
\end{proof}

\begin{proposition}[Barrier R13b: the remaining adelic gates]
\label{prop:barrier-adelic}
The single-$n$, many-place route meets the following four additional
certificates.
\begin{enumerate}
\item \emph{Partition blindness.}  The logarithm of the contribution
of primes in $(X,2X]$ is an internal partition of the fixed mass
$\log b_n$.  Heights weigh that divisor but do not locate its support.
Every algebraic arithmetization of \eqref{eq:adelic-data} pays the
$D_n^3$ tax, which is divisible at all window primes and hence has no
localizing content.
\item \emph{Unit-points versus approximation-vectors.}  The
Subspace-Theorem amplification for $S$-unit points uses coordinate
mass and form smallness on disjoint place sets.  Here the coordinate
$b_n$ itself vanishes at each moving prime, so the input is the unknown
middle-prime mass.  A manufactured second integer can help only when
the residues share a fixed lift $z'$; then every window prime divides
$n-z'$, giving only polylogarithmically many primes unless $n=z'$.
The latter is precisely the coherent apex excluded by $(\Omega)$.
\item \emph{Wrong sign and wrong disc.}  The
Andr\'e--Chudnovsky denominator theory contributes the same $D_n^3$
tax, not a denominator gain.  Available $p$-adic $G$-function
transcendence works inside a strict $p$-adic disc, whereas the Dwork
evaluations defining $Z_p$ occur at Teichm\"uller boundary points of
absolute value $1$ \cite{Andre1989,Chudnovsky1985,Dwork1969}.
Uniformly reaching that boundary would already require the missing
unit-root information.
\item \emph{Skeleton frontier and radical transparency.}  The current
frontier is summarized by
\[
\begin{array}{c|c|c}
\text{sequence}&\text{apparition skeleton}&\text{status}\\ \hline
2^n-1&\operatorname{ord}_p(2)\mid n&X^{o(1)}\text{ elementary}\\
\text{two-torus $S$-units}&\text{finite-rank lattice}&
 \text{unconditional}\\
\text{elliptic divisibility}&\text{elliptic apparition law}&
 \text{Vojta-conditional}\\
\text{Hecke coefficients}&\text{multiplicative, ensemble-level}&
 \text{open}\\
\text{Ap\'ery numbers}&\text{no known apparition law}&
 \text{beyond this frontier}.
\end{array}
\tag{8.17}\label{eq:skeleton-frontier}
\]
Moreover, an $abc$ argument based on radical versus height sees the
excess $(v_p-1)\log p$ but is transparent to the valuation-one primes
that make up the radical target $(\Omega)$.  The valuation-one phantom
therefore passes this input as well.
\end{enumerate}
\end{proposition}

\begin{proof}[Proof sketch]
The first item is invariance of projective height under repartitioning
the finite-place contribution.  The second contrasts the
finite-rank-unit mechanism behind the Corvaja--Zannier gcd theorem
\cite{BugeaudCorvajaZannier2003} with the self-referential coordinate
$b_n$; the fixed-lift assertion follows because the product of distinct
window primes dividing the nonzero integer $n-z'$ is bounded by its
size.  The third item is the location of the known $G$-function
estimates relative to the Teichm\"uller boundary.  The fourth records
the progressively weaker multiplicative skeleton; the conditional
abelian-variety frontier is represented by Silverman's Vojta analysis
\cite{Silverman2005}.  Finally, valuation-one primes do not enlarge the
difference between the logarithm of an integer and the logarithm of
its radical, which proves the stated $abc$-transparency certificate.
\end{proof}

\subsubsection{Composition and the missing comparison}

\begin{proposition}[Two unifying principles]
\label{prop:two-principles}
The six barriers compose without circularity.
\begin{enumerate}
\item \emph{Phantom-transparency:} the five ensemble families use
statistics unchanged when the true zero-section is exchanged for the
coherent projected-pencil section, whereas $(\mathrm{SA})$ is changed.
\item \emph{Critical adelic ledger:} the single-$n$ arithmetic family
does not use that exchangeability; it stops because its
algebraic-form surplus is exactly the already-spent irrationality
margin \eqref{eq:critical-ledger}, and because heights cannot locate a
divisor's support.
\end{enumerate}
Thus the first five barriers and the sixth adelic barrier are
independent certificates for the two proposed exits.
\end{proposition}

\begin{proof}[Proof sketch]
The first five barrier propositions inspect ensemble or fixed-quotient
data.  Proposition~\ref{prop:barrier-adelic} instead fixes one $n$ and
audits the product-formula ledger across its places.  Its obstruction
does not assume the ensemble phantom, while the phantom remains
compatible with its radical and height data.  No premise of either
certificate is therefore borrowed from the desired conclusion of the
other.
\end{proof}

\begin{remark}[Scope of the six-fold map]
Current multi-place tools factor through one or more of three kinds of
input: heights, separate per-place data, and ensemble statistics.
Partition blindness, the boundary-disc obstruction, and
phantom-transparency respectively block those three inputs.  In this
precise methodological sense the audit leaves no seventh current
family.  This is an optimality statement for the reduction and the
audited technology, not a theorem that an unknown kind of comparison
cannot exist.
\end{remark}

\begin{remark}[The missing product formula for periods]
\label{rem:period-product-formula}
The remaining target may be named a \emph{product formula for periods}:
a place-mixing inequality for the Ap\'ery motive at one integer index
$n$, comparing its Betti/de Rham realization, represented by the
archimedean period remainder $\ell_n$, with its crystalline
realizations at every $p\in(X,2X]$, represented by digit evaluations of
the Frobenius crystal at Teichm\"uller boundary points.  Such an
inequality would have to imply a power saving in $(\Omega)$ without
factoring through an algebraic avatar that pays the universal $D_n^3$
tax.  This is a description of the required new conjectural input, not
a theorem supplied here.  Neither the standard GRH-type information
audited above nor $abc$ provides this place-mixing comparison.
\end{remark}

\subsubsection{The seventh barrier: constructing the comparison}

Write $(\mathrm B)$ for the terminal exact-collision target
\eqref{eq:Ncoll}, or, under $(\mathrm C1)$ and after renaming positive
exponents, for either of its equivalent forms $(\mathrm{SA})$ and
$(\Omega)$.  The preceding discussion characterizes the desired
output of a place-mixing construction.  Indeed, if a nonzero integer
$C_n$ is divisible by every prime in $(X,2X]$ which divides $b_n$,
then
\[
 \omega_{(X,2X]}(b_n)\leq\frac{\log|C_n|}{\log X}.
 \tag{8.18}\label{eq:avatar-bound}
\]
Thus an avatar-integer proof of $(\Omega)$ would have to produce a
window-supported $C_n$ with
$\log|C_n|\ll X^{1-\eta}\log X$.  The ordinary holonomic avatars have
logarithmic height comparable with $n$, which is
$X(\log X)^{O(1)}$ in the in-range window.  The missing saving is a
genuine power of $X$, not a logarithmic refinement of the height
ledger.

\begin{proposition}[Barrier R14: the construction gate]
\label{prop:barrier-construction}
The three proposed constructions of the comparison in
Remark~\ref{rem:period-product-formula} meet the following five
certificates.  Each assertion is relative to the indicated
construction formalism.
\begin{enumerate}
\item \emph{Functoriality and the invariant ring.}  A bounded complexity
scalar constructed functorially from the fixed rank-$2$ difference
module lies in its invariant calculus.  Parts (NI-1) and (NI-2) of
Proposition~\ref{prop:barrier-dynamical} show that this calculus
factors through the determinant, or Casoratian, line.  Its
logarithmic height is $O(\log n)$, but its prime support is confined to
the recurrence singularities $n\equiv0,-1\pmod p$ and is therefore
window-blind.  In the Birkhoff-weight calculus, every nonzero weight is
on the exponential scale $e^{c n}$.  The resulting height spectrum is
gapped: it contains no window-supported avatar in the intermediate
slot required by \eqref{eq:avatar-bound}.  Multiplying the available
per-prime continuant certificates across $p$ simply restores the full
height ledger.

\item \emph{No cross-prime interpolation.}  The digit value
$b_r\bmod p$ is represented by a finite-field hypergeometric datum
whose parameter $r/(p-1)$ moves with $p$ along $r=n\bmod p$
\cite{KalitaChetry2017}.  These formulas therefore do not place the
moving data in the fixed-object interpolation spaces used by motivic
or $p$-adic $L$-function constructions.
The fixed-parameter channels are sparse for an elementary reason: if
\[
 \frac{r}{p-1}=\frac ak
 \quad\text{and}\quad r\equiv n\pmod p
\]
with $a/k$ fixed, then $p\mid kn+a$.  For fixed $a/k$ these channels
contain only $O(\log n)$ prime divisors, and the central channel
$a/k=1/2$ is exactly $p\mid2n+1$, where the Ahlgren--Ono modular-form
congruence applies \cite{AhlgrenOno2000}.  Existing $L$-value and
regulator constructions for a fixed object consequently see either
these torsion-digit channels or global heights; neither output locates
the generic moving digit.

\item \emph{Elliptic-divisibility calibration.}  Even an exact
place-by-place height formula is not a locating theorem.  Elliptic
divisibility sequences possess such a formula and, in addition, a
fixed-quotient apparition law, yet their effective window problem
still passes through growing-level Chebotarev estimates.  The precise
calibration is recorded in Remark~\ref{rem:EDS-calibration} below.

\item \emph{Boundary and non-overconvergence.}  Fontaine--Laffaille and
integral $p$-adic Hodge comparison theorems provide functorial
filtration and lattice divisibilities
\cite{FontaineLaffaille1982,BhattMorrowScholze2018}.  Their uniform
part is independent of the digit and hence reproduces the universal
$D_n^3$ tax.  At Teichm\"uller digit points, the available
sequence-specific output remains the Beukers--Vlasenko vertical tower
\[
 b_{mp^k}\equiv b_{mp^{k-1}}\pmod {p^{3k}}
 \qquad(p\geq5),
\]
in its applicable range \cite{Coster1988,BeukersVlasenko2021}.  This compares
$p$-power indices at one prime; it neither couples two primes nor
locates $n\bmod p$ inside $Z_p$.  Integral refinement does not lower
the intrinsic linear-in-$p$ complexity of the digit polynomial
$\Phi_p$ or supply the missing uniform continuation from a strict
$p$-adic disc to all Teichm\"uller boundary points.

\item \emph{Direction of place mixing.}  The audited local-to-global
theorems run from strong information at many primes to a rigid global
conclusion: Fontaine's nonexistence theorem, the
Grothendieck--Katz $p$-curvature paradigm, and arithmetic
algebraization are standard examples
\cite{Fontaine1985,Katz1972,Andre1989,Bost2001,Chudnovsky1985}.
Target $(\mathrm B)$ asks for the reverse direction, from soft
archimedean information to a cap on a conspiracy of self-characteristic
local events.  The Pila--Zannier strategy likewise requires a
bounded-complexity family of special loci
\cite{PilaZannier2008}; the available defining polynomials
$\Phi_p$ have degree $4(p-1)$, and no bounded-complexity replacement
has been supplied.  Consequently this strategy, without a new
uniformity theorem, re-enters the same height ledger rather than producing
\eqref{eq:avatar-bound}.  In this audited sense there is no eighth
current family beyond the construction gate.
\end{enumerate}
\end{proposition}

\begin{proof}[Proof sketch]
The first item is Proposition~\ref{prop:barrier-dynamical} applied to
functorial scalars, together with the two Birkhoff weights used in its
proof.  The invariant line gives only the Casoratian factors; a
non-invariant covariant has nonzero difference weight.  The second
item uses the displayed congruence $kn\equiv-a\pmod p$; counting prime
divisors of the fixed nonzero integer $kn+a$ gives the stated
logarithmic bound.  The third item is
Remark~\ref{rem:EDS-calibration}.  In the fourth item, filtration
divisibility is uniform in the point, whereas the displayed Dwork
tower is vertical at a fixed prime.  The final item compares the
directions and the complexity hypotheses of the cited theorems.  None
of these arguments claims logical independence from an unknown
construction; they certify the limits of the stated mechanisms.
\end{proof}

\begin{remark}[EDS calibration: a product formula that still does not locate]
\label{rem:EDS-calibration}
Let $E/\mathbb Q$ be an elliptic curve, let $P\in E(\mathbb Q)$ be
nontorsion, and write, on a minimal model,
\[
 x(mP)=\frac{A_m}{B_m^2},\qquad (A_m,B_m)=1.
\]
At every prime of good reduction, the finite local N\'eron height has
the exact denominator term
\[
 \lambda_p(mP)=v_p(B_m)\log p.
\]
After the standard corrections at the finitely many bad places, the
local heights satisfy the exact global decomposition
\[
 \sum_v\lambda_v(mP)=\widehat h(mP)=m^2\widehat h(P).
 \tag{8.19}\label{eq:EDS-height}
\]
Moreover, if $r_p$ is the order of the reduction of $P$ in
$E(\mathbb F_p)$, then
\[
 p\mid B_m\quad\Longleftrightarrow\quad r_p\mid m,
 \qquad r_p\mid\#E(\mathbb F_p).
\]
These apparition conditions are visible in finite division and Kummer
quotients \cite{SilvermanEDS2005}.  Nevertheless, the standard route
to a one-sequence window cap must count Frobenius conditions in levels
growing with the window; estimates of the needed uniform strength are
GRH-grade, while unconditional effective Chebotarev is too weak for
that route.  For two independent divisibility sequences, Silverman's
strong gcd cap used in this calibration is conditional on Vojta's
conjecture \cite{Silverman2005}.  Thus the EDS model possesses both an exact
place-by-place ledger and a fixed-quotient apparition law, but these do
not by themselves give the required unconditional locating estimate.
The product formula is the necessary \emph{format} of the proposed
comparison, not its locating content.
\end{remark}

\begin{conjecture}[Quotient-free locating for the Ap\'ery loci]
\label{conj:quotient-free-locating}
The terminal bound \eqref{eq:Ncoll} holds for the self-characteristic
loci $n\bmod p\in Z_p$.  Equivalently under $(\mathrm C1)$, after
renaming the positive exponent, either \eqref{eq:SA} or
\eqref{eq:Omega} holds.
\end{conjecture}

This is the same numerical assertion as
Conjecture~\ref{conj:Ncoll}, not an additional hypothesis.  The new
name records the mechanism that is missing: an effective
apparition/equidistribution theorem for entropy-positive $p$-local
loci which does not first map the events to a fixed auxiliary
quotient.  No theorem in the cited $p$-adic Hodge, GRH--Chebotarev,
$abc$, Vojta, or unlikely-intersection frameworks supplies this
quotient-free statement.

\begin{proposition}[The seven-fold map and three unifying principles]
\label{prop:sevenfold-map}
Within the scope stated in the barrier propositions, the analytic,
algebraic geometry, Diophantine, dynamical, additive combinatorics,
adelic, and motivic construction barriers compose without circularity.
Their obstructions are organized by three principles:
\begin{enumerate}
\item \emph{Phantom-transparency} for the five ensemble families;
\item the \emph{critical adelic ledger} for algebraic single-index
avatars; and
\item \emph{quotient-free locating}: every locating mechanism in the
calibrated siblings factors through a fixed quotient, while the
Ap\'ery digit event has no bounded-complexity quotient in the audited
invariant calculus.
\end{enumerate}
The missing statement after all seven barriers is
Conjecture~\ref{conj:quotient-free-locating}, not merely the existence
of a period comparison.
\end{proposition}

\begin{proof}[Proof sketch]
Proposition~\ref{prop:two-principles} composes the first six barriers.
Proposition~\ref{prop:barrier-construction} does not use a replacement
zero-section or spend a second copy of the adelic surplus: it audits
the proposed construction of the comparison itself.  The EDS
calibration then separates a place-by-place ledger from the additional
act of locating its support.  Hence the seventh certificate adds the
third principle without borrowing the desired collision bound.
\end{proof}

\begin{remark}[Relation to the six-fold scope statement]
The phrase ``no seventh current family'' above referred to output
formats already covered by heights, separate per-place data, and
ensemble statistics.  Proposition~\ref{prop:barrier-construction}
instead treats the previously unaudited attempt to \emph{construct}
the proposed comparison and records it as the seventh barrier.
Its direction-of-place-mixing item preempts an eighth family assembled
from the presently cited tools.
Both scope claims remain methodological, not logical-independence
theorems.
\end{remark}

\begin{remark}[Numerical evidence against a disproof]
Exact integer/rational computation through $n=6000$ gives
\[
 \max_{n\leq N}\log G_n\approx0.75(\log N)^2;
\]
over three decades the observed ratio to $(\log N)^2$ stays between
about $0.6$ and $0.9$.  The worst cases are diffuse rather than
single-prime spikes.  At the recorded worst indices the largest prime
factor always occurs with $v_p=1$, and no in-range full-multiplicity
phantom occurs.  This is the Poisson-independence scale predicted by
the zero-count data and makes the available disproof mechanisms
numerically nonviable.  It is not a proof of the pointwise conjecture:
the finite computation cannot exclude an extremely sparse later
phase transition.
\end{remark}

\section{Conclusion}

The density form of Problem~3.2 is unconditional.  The reflection
functional equation, the gap-eliminant theorem for $U_g$, the cross-arc
lemma, and the structured-resonance theorem are unconditional as well.
The only delicate dyadic blocks are in range, and there the reduction is
\[
 \begin{aligned}
 \text{terminal exact-collision conjecture}
 &\Longrightarrow (P_2)^G\\
 &\Longrightarrow \sup K_X\text{ has a power saving}\\
 &\Longrightarrow \text{pointwise Problem~3.2}.
 \end{aligned}
\]
Under the mass conjecture $(\mathrm C1)$, the terminal conjecture is
equivalent, after renaming positive exponents, to the uniform
single-arc cap $(\mathrm{SA})$, the middle-prime cap $(\Omega)$, and the
power-saving supremum target.  The present argument does not claim that
the qualitative conclusion $\log G_n=o(n)$ conversely forces a fixed
power saving.

The terminal problem is a holonomic shifted-prime divisor problem in
the circle of Hooley and Murty--Rosen--Silverman, sharpened to exact
collisions.  Lemma~\ref{lem:Mersenne} shows how a compressive
multiplicative invariant solves the Mersenne analogue; no counterpart
is known for Ap\'ery zeros.  Phantom locality, the bandwidth model
R9.1, and the Lang--Trotter implication R9.2 certify the limits of the
currently banked divisibility and equidistribution routes.  What
remains is genuinely scale zero: either find a cross-prime compressive
invariant, or prove an exact-collision theorem for holonomic values.

The six-fold map sharpens this status.  The density theorem is fully
unconditional.  At the fixed-power endpoint, and under the mass
conjecture $(\mathrm C1)$ where required, $(\mathrm{SA})$,
$(\Omega)$, and $(\mathrm{NC})$ are equivalent after exponent renaming,
and each implies pointwise Problem~3.2.  This is the named
Murty--Rosen--Silverman holonomic-sibling exact-collision problem.  The
computations in the appendix are compatible with the Poisson model and
leave a large numerical margin.  As stated above, no converse from the
merely qualitative conclusion to a fixed power saving is claimed.

The reduction is optimal relative to the audited technology: five
ensemble families---analytic, algebraic-geometric, Diophantine,
dynamical, and additive combinatorics---share the phantom-transparency
principle.  The single-index adelic family instead stops at the critical
ledger.  The two certificates compose without circularity.
Heights, separate per-place information, and ensemble statistics
exhaust the input formats of the present multi-place tools, in the
qualified sense of Remark~\ref{rem:period-product-formula}.

A fully unconditional pointwise theorem therefore needs genuinely new
place-mixing input of the type described there: a product formula for
periods comparing the Betti/de Rham and crystalline realizations of the
Ap\'ery motive at one integer point.  Standard GRH-type hypotheses do
not formulate this comparison, and the valuation-one radical target is
transparent to $abc$.  The paper thus resolves the status of the
pointwise problem while leaving its named scale-zero conjecture open.

The construction audit in Proposition~\ref{prop:barrier-construction}
refines this last description.  The elliptic-divisibility calibration
shows that a product formula for periods would provide the right
place-by-place \emph{format}, but would not by itself locate the
window-prime support.  The exact fixed-power reduction is
\[
 \text{pointwise Problem~3.2}
 \quad\Longleftarrow\quad
 (\mathrm{SA})\Longleftrightarrow(\Omega)
 \Longleftrightarrow(\mathrm{NC}),
\]
where the equivalences use $(\mathrm C1)$ when required and allow
renaming the saving exponent.  The reverse implication from the
qualitative pointwise statement is not claimed.

The optimality certificate is therefore seven-fold: the five ensemble
families, the single-index adelic family, and the
motivic-comparison/construction gate R14.  It is summarized by
phantom-transparency, the critical adelic ledger, and quotient-free
locating.  Exact computation through $n=6000$ supports, over the tested
range, $\max_{n\leq N}\log G_n\approx0.75(\log N)^2$ and exhibits
no full-multiplicity phantom, but remains numerical evidence only.
The missing mathematics is
Conjecture~\ref{conj:quotient-free-locating}: an effective locating
theorem for entropy-positive self-characteristic local loci without a
fixed quotient.  None of the audited consequences of $p$-adic Hodge
theory, GRH, $abc$, Vojta's conjecture, or unlikely-intersection
technology supplies it.  In this precise sense the required theorem is
new, rather than a new combination of the seven audited families.

\appendix

\section{Verification summary}

\begin{itemize}
\item Lemmas~\ref{lem:no-consecutive}--\ref{lem:zero-fiber},
Lemma~\ref{lem:block-law}, and Theorem~\ref{thm:support-law} were checked
exactly on all $8434$ triples $(p,q,r)$ with $7\leq p\leq200$, including
the support equivalence against exact $G_n$.  The zero-pair certificate
was checked on all $232$ consecutive zero pairs with $p<3000$, and the
evaluation identities and reflection functional equation were checked
exactly through $r=9$.

\item Lemma~\ref{lem:alpha-six} was checked for all $31$ primes
$7\leq p\leq139$.

\item Lemma~\ref{lem:boundary} was checked for $q=2,3$ and all
$p\leq139$, in addition to its role in the $8434$-triple support check.

\item After noncentral roots are grouped into reflection pairs, the
$Z_b$ data are close to a Poisson law of mean $1/2$; the maximum
$|Z_p|$ is $16$ among the $2704$ primes up to $80000$.  The empirical
mean of $|Z_p|$ is about $1$; the value $1/8$ belongs to a different
collision statistic and is not a zero-count mean.  Thus the $p^{2/3}$
bound is far from tight.

\item In the pointwise probe,
\[
 B(n)=\#\{\text{bad primes }p>\sqrt n\}\leq5
\]
for every $n\leq10^4$.

\item At five dyadic scales, the all-range and in-range data were
\[
 \begin{array}{c|ccccc}
 X&400&800&1600&3200&6400\\ \hline
 \max_{m<X^2}K_X(m)&3&4&4&4&5\\
 P_{2,X}&6&28&49&109&276\\
 \max_{y\in[2X,X\log X]}K_X(y)&2&2&3&3&3
 \end{array}
\]
These are numerical checks only.

\item The reflection functional equation had no failure for
$p\leq12809$; its termwise binomial proof was checked for $p\leq809$.

\item For $n=250,500,1000,2000,3000,4000$, respectively, the numbers of
distinct prime divisors of $b_n$ in $(\sqrt n,n]$ were
\[
 2,3,0,0,0,2,
\]
and the corresponding full-multiplicity logarithmic ratios were
\[
 0.030,0.029,0,0,0,0.003.
\]
The largest observed middle-prime
multiplicity was $1$.

\item The eliminant identity was checked in $116$ cases over primes from
$101$ through $1601$, with no failure.  For $g=2,3,4,5$, symbolic
calculation gave degrees $3,6,9,12$ and coefficient heights
$231$, $2.6\cdot10^5$, $6.6\cdot10^8$, and $3.6\cdot10^{12}$.

\item The $\zeta(2)$ central congruence had no mismatch for primes
$5\leq p\leq1200$.  At $X=1600$, its out-of-range count was
$C_{X,0}=3934$, while six sampled in-range counts were
$35,26,34,38,40,32$, on the random-model pair scale.  These values
support only the range correction, not an anti-collision theorem.

\item Exact bignum computation of $G_n$ through $n=6000$ found, over
three decades,
\[
 0.6\lesssim
 \frac{\max_{n\leq N}\log G_n}{(\log N)^2}
 \lesssim0.9.
\]
At the recorded worst-ratio indices
$n=523, 1667, 2098, 3471, 4855$, the gcd mass was diffuse and the
largest prime factor had exponent $1$; no in-range full-multiplicity
phantom was observed.  These are numerical checks only.
\end{itemize}

\section{v3 changelog}

\begin{itemize}
\item Replaced the obsolete depth-$3$, BAC, singleton-row, and interval-
overlap endpoints by the in-range depth-$2$ pair gate and the terminal
single-arc exact-collision conjecture.
\item Added exceptional-prime tolerance, the precise arc digit
restriction, the $(\mathrm{SA})$--$(\Omega)$ comparison, and the exact
identity \eqref{eq:Ncoll-identity}.  Power-saving equivalences now
explicitly allow exponent renaming.
\item Added the cross-arc lemma, the structured-resonance theorem, and
the Ap\'ery--Hasse and projective-Frobenius descriptions of $Z_p$.
\item Added the Mersenne fibration calibration and barriers R9.1 and
R9.2.  R9.1 is labelled as a model because the source notes do not give
an actual root-set realization.
\item Corrected the zero-count mean to approximately $1$, retained the
Rhin--Viola constant $5.24673\ldots$, restricted Lemma~\ref{lem:Mersenne}
to the admissible window, and removed the false fixed-exponent and
exceptional-$y$ equivalences.
\end{itemize}

\section{v4 changelog}

\noindent Added the sieve-complexity invariant, the complete six-fold barrier map and adelic ledger, the product-formula-for-periods target, and the corresponding pointwise-status conclusion.

\section{v5 changelog}

\noindent Added the R14 construction gate, EDS calibration, quotient-free-locating frontier, seven-fold barrier map, disproof probe, and refined pointwise-status conclusion.

\begin{thebibliography}{99}
\raggedright

\bibitem{AhlgrenOno2000}
Scott Ahlgren and Ken Ono,
\emph{A Gaussian hypergeometric series evaluation and Ap\'ery number
congruences},
J. Reine Angew. Math. \textbf{518} (2000), 187--212.
DOI: \texttt{10.1515/crll.2000.004}.

\bibitem{Andre1989}
Yves Andr\'e,
\emph{$G$-Functions and Geometry},
Aspects of Mathematics, vol.~E13, Vieweg, Braunschweig, 1989.
DOI: \texttt{10.1007/978-3-663-14108-2}.

\bibitem{Apery1979}
Roger Ap\'ery,
\emph{Irrationalit\'e de $\zeta(2)$ et $\zeta(3)$},
Ast\'erisque \textbf{61} (1979), 11--13.

\bibitem{Beukers1985}
Frits Beukers,
\emph{Some congruences for the Ap\'ery numbers},
J. Number Theory \textbf{21} (1985), no.~2, 141--155.
DOI: \texttt{10.1016/0022-314X(85)90047-2}.

\bibitem{BeukersVlasenko2021}
Frits Beukers and Masha Vlasenko,
\emph{Dwork crystals II},
Int. Math. Res. Not. IMRN \textbf{2021} (2021), no.~6, 4427--4444.
DOI: \texttt{10.1093/imrn/rnaa120}.

\bibitem{BellonViallet1999}
Marc P. Bellon and Claude-Michel Viallet,
\emph{Algebraic entropy},
Comm. Math. Phys. \textbf{204} (1999), 425--437.
DOI: \texttt{10.1007/s002200050652}.

\bibitem{BettinChandee2018}
Sandro Bettin and Vorrapan Chandee,
\emph{Trilinear forms with Kloosterman fractions},
Adv. Math. \textbf{328} (2018), 1234--1262.
DOI: \texttt{10.1016/j.aim.2018.01.026}.

\bibitem{BhattMorrowScholze2018}
Bhargav Bhatt, Matthew Morrow, and Peter Scholze,
\emph{Integral $p$-adic Hodge theory},
Publ. Math. Inst. Hautes \'Etudes Sci. \textbf{128} (2018), 219--397.
DOI: \texttt{10.1007/s10240-019-00102-z}.

\bibitem{BombieriFriedlanderIwaniec1986}
Enrico Bombieri, John B. Friedlander, and Henryk Iwaniec,
\emph{Primes in arithmetic progressions to large moduli},
Acta Math. \textbf{156} (1986), 203--251.
DOI: \texttt{10.1007/BF02399204}.

\bibitem{Bost2001}
Jean-Beno\^\i t Bost,
\emph{Algebraic leaves of algebraic foliations over number fields},
Publ. Math. Inst. Hautes \'Etudes Sci. \textbf{93} (2001), 161--221.
DOI: \texttt{10.1007/s10240-001-8191-3}.

\bibitem{BugeaudCorvajaZannier2003}
Yann Bugeaud, Pietro Corvaja, and Umberto Zannier,
\emph{An upper bound for the G.C.D. of $a^n-1$ and $b^n-1$},
Math. Z. \textbf{243} (2003), 79--84.

\bibitem{CFVMZ2026}
Xavier Caruso, Florian F\"urnsinn, Daniel Vargas-Montoya, and
Wadim Zudilin,
\emph{Galois groups of Ap\'ery-like series modulo primes},
arXiv:2510.23298 (2025; revised 2026).

\bibitem{Chudnovsky1985}
David V. Chudnovsky and Gregory V. Chudnovsky,
\emph{Applications of Pad\'e approximations to Diophantine
inequalities in values of $G$-functions},
in \emph{Number Theory}, Lecture Notes in Mathematics, vol.~1135,
Springer, Berlin, 1985, 9--51.
DOI: \texttt{10.1007/BFb0074600}.

\bibitem{Coster1988}
Matthijs J. Coster,
\emph{Supercongruences},
Ph.D. thesis, Universiteit Leiden, 1988.

\bibitem{CorvajaZannier2005}
Pietro Corvaja and Umberto Zannier,
\emph{A lower bound for the height of a rational function at
$S$-unit points},
Monatsh. Math. \textbf{144} (2005), no.~3, 203--224.

\bibitem{Delaygue2018}
\'Eric Delaygue,
\emph{Arithmetic properties of Ap\'ery-like numbers},
Compos. Math. \textbf{154} (2018), no.~2, 249--274.
DOI: \texttt{10.1112/S0010437X17007552}.

\bibitem{DFI1995}
William Duke, John B. Friedlander, and Henryk Iwaniec,
\emph{Equidistribution of roots of a quadratic congruence to prime
moduli},
Ann. of Math. (2) \textbf{141} (1995), no.~2, 423--441.
DOI: \texttt{10.2307/2118527}.

\bibitem{DFI1997}
William Duke, John B. Friedlander, and Henryk Iwaniec,
\emph{Bilinear forms with Kloosterman fractions},
Invent. Math. \textbf{128} (1997), no.~1, 23--43.
DOI: \texttt{10.1007/s002220050135}.

\bibitem{Dwork1969}
Bernard Dwork,
\emph{$p$-adic cycles},
Publ. Math. Inst. Hautes \'Etudes Sci. \textbf{37} (1969), 27--115.
DOI: \texttt{10.1007/BF02684886}.

\bibitem{EvertseSchlickeweiSchmidt2002}
Jan-Hendrik Evertse, Hans Peter Schlickewei, and Wolfgang M. Schmidt,
\emph{Linear equations in variables which lie in a multiplicative
group},
Ann. of Math. (2) \textbf{155} (2002), no.~3, 807--836.
DOI: \texttt{10.2307/3062133}.

\bibitem{Fontaine1985}
Jean-Marc Fontaine,
\emph{Il n'y a pas de vari\'et\'e ab\'elienne sur $\mathbb Z$},
Invent. Math. \textbf{81} (1985), no.~3, 515--538.

\bibitem{FontaineLaffaille1982}
Jean-Marc Fontaine and Guy Laffaille,
\emph{Construction de repr\'esentations $p$-adiques},
Ann. Sci. \'Ecole Norm. Sup. (4) \textbf{15} (1982), no.~4, 547--608.
DOI: \texttt{10.24033/asens.1437}.

\bibitem{FouvryIwaniec1983}
\'Etienne Fouvry and Henryk Iwaniec,
\emph{Primes in arithmetic progressions},
Acta Arith. \textbf{42} (1983), 197--218.
DOI: \texttt{10.4064/aa-42-2-197-218}.

\bibitem{Gessel1982}
Ira Gessel,
\emph{Some congruences for Ap\'ery numbers},
J. Number Theory \textbf{14} (1982), no.~3, 362--368.
DOI: \texttt{10.1016/0022-314X(82)90071-3}.

\bibitem{Greene1987}
John Greene,
\emph{Hypergeometric functions over finite fields},
Trans. Amer. Math. Soc. \textbf{301} (1987), no.~1, 77--101.
DOI: \texttt{10.1090/S0002-9947-1987-0879564-8}.

\bibitem{HendriksSinger1999}
Peter A. Hendriks and Michael F. Singer,
\emph{Solving difference equations in finite terms},
J. Symbolic Comput. \textbf{27} (1999), no.~3, 239--259.

\bibitem{Hooley1964}
Christopher Hooley,
\emph{On the distribution of the roots of polynomial congruences},
Mathematika \textbf{11} (1964), 39--49.
DOI: \texttt{10.1112/S0025579300003466}.

\bibitem{Hooley1973}
Christopher Hooley,
\emph{On the largest prime factor of $p+a$},
Mathematika \textbf{20} (1973), no.~2, 135--143.
DOI: \texttt{10.1112/S002557930000471X}.

\bibitem{KankiMadaTokihiro2013}
Masataka Kanki, Jun Mada, and Tetsuji Tokihiro,
\emph{The space of initial conditions and the property of an almost
good reduction in discrete Painlev\'e II equations over finite fields},
J. Nonlinear Math. Phys. \textbf{20} (2013), suppl.~1, 101--109.
DOI: \texttt{10.1080/14029251.2013.862437}.

\bibitem{KalitaChetry2017}
Gautam Kalita and Arjun Singh Chetry,
\emph{Congruences for generalized Ap\'ery numbers and Gaussian
hypergeometric series},
Res. Number Theory \textbf{3} (2017), Paper No.~5.
DOI: \texttt{10.1007/s40993-016-0069-z}.

\bibitem{Katz1972}
Nicholas M. Katz,
\emph{Algebraic solutions of differential equations
($p$-curvature and the Hodge filtration)},
Invent. Math. \textbf{18} (1972), no.~1--2, 1--118.
DOI: \texttt{10.1007/BF01389714}.

\bibitem{Kilbourn2006}
Timothy Kilbourn,
\emph{An extension of the Ap\'ery number supercongruence},
Acta Arith. \textbf{123} (2006), no.~4, 335--348.
DOI: \texttt{10.4064/aa123-4-3}.

\bibitem{LucaShparlinski2008}
Florian Luca and Igor E. Shparlinski,
\emph{Arithmetic properties of Ap\'ery numbers},
J. Lond. Math. Soc. (2) \textbf{78} (2008), no.~3, 545--562.
DOI: \texttt{10.1112/jlms/jdn031}.

\bibitem{MalikStraub2016}
Amita Malik and Armin Straub,
\emph{Divisibility properties of sporadic Ap\'ery-like numbers},
Res. Number Theory \textbf{2} (2016), no.~1, Paper No.~5, 26~pp.
DOI: \texttt{10.1007/\allowbreak s40993-\allowbreak016-\allowbreak0036-8}.

\bibitem{McCarthy2012}
Dermot McCarthy,
\emph{Extending Gaussian hypergeometric series to the $p$-adic
setting},
Int. J. Number Theory \textbf{8} (2012), no.~7, 1581--1612.
DOI: \texttt{10.1142/S1793042112500844}.

\bibitem{Mortenson2005}
Eric Mortenson,
\emph{Supercongruences for truncated ${}_{n+1}F_n$ hypergeometric series
with applications to certain weight three newforms},
Proc. Amer. Math. Soc. \textbf{133} (2005), no.~2, 321--330.
DOI: \texttt{10.1090/S0002-9939-04-07697-X}.

\bibitem{MurtyRosenSilverman1996}
M. Ram Murty, Michael Rosen, and Joseph H. Silverman,
\emph{Variations on a theme of Romanoff},
Internat. J. Math. \textbf{7} (1996), no.~3, 373--391.
DOI: \texttt{10.1142/\allowbreak S0129167X\allowbreak96000220}.

\bibitem{NewtonThorne2021}
James Newton and Jack A. Thorne,
\emph{Symmetric power functoriality for holomorphic modular forms},
Publ. Math. Inst. Hautes \'Etudes Sci. \textbf{134} (2021), 1--116.

\bibitem{PilaZannier2008}
Jonathan Pila and Umberto Zannier,
\emph{Rational points in periodic analytic sets and the
Manin--Mumford conjecture},
Atti Accad. Naz. Lincei Cl. Sci. Fis. Mat. Natur. \textbf{19} (2008),
no.~2, 149--162.
DOI: \texttt{10.4171/RLM/514}.

\bibitem{RhinViola2001}
Georges Rhin and Carlo Viola,
\emph{The group structure for $\zeta(3)$},
Acta Arith. \textbf{97} (2001), no.~3, 269--293.
DOI: \texttt{10.4064/aa97-3-6}.

\bibitem{Romanoff1934}
N. P. Romanoff,
\emph{\"Uber einige S\"atze der additiven Zahlentheorie},
Math. Ann. \textbf{109} (1934), no.~1, 668--678.
DOI: \texttt{10.1007/BF01449161}.

\bibitem{RowlandYassawiKrattenthaler2021}
Eric Rowland, Reem Yassawi, and Christian Krattenthaler,
\emph{Lucas congruences for the Ap\'ery numbers modulo $p^2$},
Integers \textbf{21} (2021), Paper No.~A20, 15~pp.

\bibitem{Schmidt1980}
Wolfgang M. Schmidt,
\emph{Diophantine Approximation},
Lecture Notes in Mathematics, vol.~785, Springer, Berlin, 1980.

\bibitem{Schuett2009}
Matthias Sch\"utt,
\emph{CM newforms with rational coefficients},
Ramanujan J. \textbf{19} (2009), no.~2, 187--205.
DOI: \texttt{10.1007/s11139-008-9147-8}.

\bibitem{SilvermanEDS2005}
Joseph H. Silverman,
\emph{$p$-adic properties of division polynomials and elliptic
divisibility sequences},
Math. Ann. \textbf{332} (2005), no.~2, 443--471.
DOI: \texttt{10.1007/s00208-004-0608-0}.

\bibitem{Silverman2005}
Joseph H. Silverman,
\emph{Generalized greatest common divisors, divisibility sequences,
and Vojta's conjecture for blowups},
Monatsh. Math. \textbf{145} (2005), no.~4, 333--350.

\bibitem{StienstraBeukers1985}
Jan Stienstra and Frits Beukers,
\emph{On the Picard--Fuchs equation and the formal Brauer group of
certain elliptic K3-surfaces},
Math. Ann. \textbf{271} (1985), no.~2, 269--304.
DOI: \texttt{10.1007/BF01455990}.

\bibitem{Stewart2012}
Cameron L. Stewart,
\emph{On divisors of Lucas and Lehmer numbers},
Acta Math. \textbf{211} (2013), no.~2, 291--314.

\bibitem{Sun2012}
Zhi-Wei Sun,
\emph{On sums of Ap\'ery polynomials and related congruences},
J. Number Theory \textbf{132} (2012), no.~11, 2673--2699.
DOI: \texttt{10.1016/j.jnt.2012.05.014}.

\bibitem{WangSun2024}
Chen Wang and Zhi-Wei Sun,
\emph{$p$-adic analogues of hypergeometric identities and their
applications},
Nanjing Univ. J. Math. Biquarterly \textbf{41} (2024), no.~1, 34--56.
arXiv:1910.06856.

\end{thebibliography}

\end{document}
